% ============================================================
%  CHAPTER I -- The Fourier Transform   (printed pp. 1-36)
%  section 1 done by hand as the reference model; sections 2-4 follow.
% ============================================================
\chapter{The Fourier Transform}

In this chapter we introduce the Fourier transform and study its more
elementary properties. Since most of the material of this chapter is rather
standard our treatment here will be brief. We begin by considering the behavior
of the Fourier transform on the spaces $L^{1}(E_{n})$ and $L^{2}(E_{n})$. This
will be done in the first two sections. The Fourier transform's formal aspects
are more easily described in the context of distributions; therefore, in the
third section we extend its definition to the space of tempered distributions.
The reader will note that in this chapter we are mainly exploiting the
translation structure of Euclidean\index{norm@norm!Euclidean@Euclidean} spaces. In the following chapters (and
specifically in Chapter IV), however, the action of rotations on these spaces
plays an important role.

\section{The Basic \texorpdfstring{$L^{1}$}{L1} Theory of the Fourier Transform}

We begin by introducing some notation that will be used throughout this work.
$E_{n}$ denotes $n$-dimensional (real) Euclidean space. We consistently write
$x = (x_{1}, x_{2}, \ldots, x_{n})$, $y = (y_{1}, y_{2}, \ldots, y_{n})$,
$\ldots$ for the elements of $E_{n}$. The inner product\index{inner product@inner product} of $x, y \in E_{n}$ is
the number $x \cdot y = \sum_{j=1}^{n} x_{j} y_{j}$; the norm of $x \in E_{n}$
is the (nonnegative) number $|x| = \sqrt{x \cdot x}$; furthermore,
$dx = dx_{1} dx_{2} \cdots dx_{n}$ denotes the element of ordinary Lebesgue
measure.

We will deal with various spaces of functions defined on\index{norm@norm!Lp@$L^{p}$} $E_{n}$. The simplest
of these are the $L^{p} = L^{p}(E_{n})$ spaces, $1 \leq p < \infty$, of all
measurable functions $f$ such that
$\|f\|_{p} = (\int_{E_{n}} |f(x)|^{p} \, dx)^{1/p} < \infty$. The number
$\|f\|_{p}$ is called the $L^{p}$ norm of $f$. The space $L^{\infty}(E_{n})$
consists of all essentially bounded functions on $E_{n}$ and, for
$f \in L^{\infty}(E_{n})$, we let $\|f\|_{\infty}$ be the essential supremum of
$|f(x)|$, $x \in E_{n}$.%
\footnote{We say $f$ is equivalent to $g$ if $f(x) = g(x)$ for almost every
$x \in E_{n}$, whenever $f \in L^{p}(E_{n})$, $1 \leq p \leq \infty$. If we
consider the equivalence classes obtained from this relation and define the
$norm$ of a class to be the norm of any one of its representatives (clearly
$\|f\|_{p} = \|g\|_{p}$ if $f$ is equivalent to $g$) we obtain a Banach space.
We shall denote this space by $L^{p}(E_{n})$ as well. It will be obvious from
the context which of these spaces $L^{p}(E_{n})$ is under discussion.}
Often, the space $C_{0}$ of all continuous functions vanishing at infinity,
with the $L^{\infty}$ norm just described, arises more naturally than
$L^{\infty} = L^{\infty}(E_{n})$. Unless otherwise specified, all functions are
assumed to be complex valued; it will be assumed, throughout the book, that all
functions are (Borel) measurable.

If $f \in L^{1}(E_{n})$ the Fourier transform\index{Fourier transform@Fourier transform} of $f$ is the function $\hat{f}$
defined by letting
\[
\hat{f}(x)=\int_{E_{n}}f(t)e^{-2\pi i x\cdot t}\,dt
\]
for all $x \in E_{n}$. It is easy to establish the following results:

\begin{theorem}\label{thm:c1-1-1}
\textup{(a)} The mapping $f \to \hat{f}$ is a bounded linear transformation
from $L^{1}(E_{n})$ into $L^{\infty}(E_{n})$. In fact
$\|\hat{f}\|_{\infty} \leq \|f\|_{1}$;

\textup{(b)} If $f \in L^{1}(E_{n})$ then $\hat{f}$ is uniformly continuous.
\end{theorem}

\begin{theorem}[Riemann-Lebesgue]\label{thm:c1-1-2}\index{Riemann-Lebesgue theorem@Riemann-Lebesgue theorem}
If $f \in L^{1}(E_{n})$ then $\hat{f}(x) \to 0$ as $|x| \to \infty$; thus, in
view of the last result, we can conclude that $\hat{f}$ belongs to the class
$C_{0}$.
\end{theorem}

Theorem~\ref{thm:c1-1-1} is obvious; moreover, so is
Theorem~\ref{thm:c1-1-2} when $f$ is the characteristic function of the
$n$-dimensional interval
$I = \{x \in E_{n}; a_{1} \leq x_{1} \leq b_{1}, \ldots, a_{n} \leq x_{n} \leq b_{n}\}$
(for we can calculate $\hat{f}$ explicitly as an iterated integral). The same is
therefore true for a finite linear combination of such characteristic functions.
The result for a general $f \in L^{1}(E_{n})$ follows easily by approximating
$f$ in the $L^{1}$ norm by such a linear combination $g$; for then
$f = g + (f - g)$, where $\hat{f} - \hat{g}$ is uniformly small (by
Theorem~\ref{thm:c1-1-1}, part (a)) while $\hat{g}(x) \to 0$ as
$|x| \to \infty$.

Theorem~\ref{thm:c1-1-2} gives a necessary condition for a function to be a
Fourier transform. Belonging to the class $C_{0}$, however, is far from being
sufficient (see \eqref{eq:c1-4-1}). There seems to be no simple satisfactory
condition characterizing Fourier transforms of functions in $L^{1}(E_{n})$.

The above definition of the Fourier transform extends immediately to finite
Borel measures: if $\mu$ is such a measure on $E_{n}$ we define $\hat{\mu}$ by
letting
\[
\hat{\mu}(x)=\int_{E_{n}}e^{-2\pi i x\cdot t}\,d\mu(t).
\]
Theorem~\ref{thm:c1-1-1} is valid for these Fourier transforms if we replace the
$L^{1}$-norm by the total variation of $\mu$.

In addition to the vector-space operations, $L^{1}(E_{n})$ is endowed with a
``multiplication'' making this space a Banach algebra. This operation, called
convolution\index{convolution@convolution}, is defined in the following way: If $f$ and $g$ belong to
$L^{1}(E_{n})$ their convolution $h = f * g$ is the function whose value at
$x \in E_{n}$ is
\[
h(x)=\int_{E_{n}}f(x-y)g(y)\,dy.
\]

One can show by an elementary argument that $f(x - y)g(y)$ is a measurable
function of the two variables $x$ and $y$. It then follows immediately from
Fubini's theorem on the interchange of the order of integration that
$h \in L^{1}(E_{n})$ and $\|h\|_{1} \leq \|f\|_{1}\|g\|_{1}$. Furthermore, this
operation is commutative and associative. More generally, $h = f * g$ is
defined whenever $f \in L^{p}(E_{n})$, $1 \leq p \leq \infty$, and
$g \in L^{1}(E_{n})$. In fact we have the following result:

\begin{theorem}\label{thm:c1-1-3}
If $f \in L^{p}(E_{n})$, $1 \leq p \leq \infty$, and $g \in L^{1}(E_{n})$ then
$h = f * g$ is well defined and belongs to $L^{p}(E_{n})$. Moreover,
\[
\|h\|_{p}\leq\|f\|_{p}\|g\|_{1}.
\]
\end{theorem}

It is clear that
$|h(x)| \leq \int_{E_{n}} |f(x - y)| \, |g(y)| \, dy$. Thus the desired result
is an easy consequence of Minkowski's integral inequality:
\begin{align*}
\left(\int_{E_{n}}|h(x)|^{p}\,dx\right)^{1/p}
&\leq\int_{E_{n}}\left(\int_{E_{n}}|f(x-y)|^{p}\,dx\right)^{1/p}|g(y)|\,dy\\
&=\|f\|_{p}\|g\|_{1}.
\end{align*}

As was the case for the Fourier transform, we can extend this operation to
include finite Borel measures: if $\mu$ is such a measure on $E_{n}$ we define
$h = f * d\mu$ by letting
\[
h(x)=\int_{E_{n}}f(x-y)\,d\mu(y),
\]
for $x \in E_{n}$ and $f \in L^{p}$. Theorem~\ref{thm:c1-1-3} is valid for these
convolutions if we replace the $L^{1}$-norm of $g$ by the total variation of
$\mu$.

An essential feature of harmonic analysis is the fact that the Fourier
transform of the convolution of two functions is the (pointwise) product of
their Fourier transforms. More precisely, the following result is an easy
consequence of the definitions:

\begin{theorem}\label{thm:c1-1-4}
If $f$ and $g$ belong to $L^{1}(E_{n})$ then
\[
(f*g)^{\wedge}=\hat{f}\hat{g}.%
\footnote{We shall use this notation consistently: $(\ldots)^{\wedge}$ denotes
the Fourier transform of $(\ldots)$.}
\]
\end{theorem}

Many other important operations of analysis have particularly simple relations
with the Fourier transform. For example, if we let $\tau_{h}$ denote translation\index{translation@translation}
by $h \in E_{n}$ (by this we mean the operator mapping the function $g(x)$ into
the function $g(x - h)$) we have
\begin{equation}\label{eq:c1-1-5}
\begin{aligned}
\textup{(i)}\quad &(\tau_{h}f)^{\wedge}(x)=e^{-2\pi i h\cdot x}\hat{f}(x),\\
\textup{(ii)}\quad &(e^{2\pi i t\cdot h}f(t))^{\wedge}(x)=(\tau_{h}\hat{f})(x),
\end{aligned}
\end{equation}
whenever $f \in L^{1}(E_{n})$.

If $a > 0$ we let $\delta_{a}$ denote dilation\index{dilation@dilation} by $a$; that is, $\delta_{a}$ is
the operator mapping the function $g(x)$ into the function $g(ax)$. Whenever
$f \in L^{1}(E_{n})$ we then have
\begin{equation}\label{eq:c1-1-6}
a^{n}(\delta_{a}f)^{\wedge}(x)=\hat{f}(a^{-1}x).
\end{equation}

Differentiation and Fourier transformation are related in the following way:

\begin{theorem}\label{thm:c1-1-7}
Suppose $f \in L^{1}(E_{n})$ and $x_{k} f(x) \in L^{1}(E_{n})$, where $x_{k}$ is
the $k$-th coordinate function. Then $\hat{f}$ is differentiable\index{differentiable@differentiable!in Lp norm@in $L^{p}$ norm} with respect to
$x_{k}$ and
\[
\frac{\partial\hat{f}}{\partial x_{k}}(x)=(-2\pi i t_{k}f(t))^{\wedge}(x).
\]
\end{theorem}

\begin{proof}
Letting $h = (0, \ldots, h_{k}, \ldots, 0)$ be a nonzero vector along the
$k$-th coordinate axis, we have, by part (ii) of \eqref{eq:c1-1-5} and the
Lebesgue dominated convergence theorem,
\[
\frac{\hat{f}(x+h)-\hat{f}(x)}{h_{k}}
=\left\{\left(\frac{e^{-2\pi i t\cdot h}-1}{h_{k}}\right)f(t)\right\}^{\wedge}(x)
\to(-2\pi i t_{k}f(t))^{\wedge}(x)
\]
as $h_{k} \to 0$.
\end{proof}

Theorem~\ref{thm:c1-1-7} asserts that applying the Fourier transform after
multiplying by the $k$-th coordinate function is equivalent (up to a
multiplicative constant) to taking the partial derivative with respect to the
$k$-th variable of the Fourier transform. It is also true that the Fourier
transforms of such partial derivatives are obtainable (again, up to a
multiplicative constant) by multiplying the Fourier transform by the
corresponding coordinate functions. We shall encounter many versions of this
result. In order to make a precise statement of one of these versions (perhaps
the simplest to prove) we introduce the following concept: We say that $f$ is
differentiable in the $L^{p}$ norm with respect to $x_{k}$ whenever
$f \in L^{p}(E_{n})$ and there exists a $g$ in $L^{p}(E_{n})$ such that
\[
\left(\int_{E_{n}}\left|\frac{f(x+h)-f(x)}{h_{k}}-g(x)\right|^{p}dx\right)^{1/p}\to0,
\]
as $h_{k} \to 0$ (we are using the notation established in the proof of
Theorem~\ref{thm:c1-1-7}). The function $g$ is called the partial derivative of
$f$ (with respect to $x_{k}$) in the $L^{p}$ norm.

Applying part (i) of \eqref{eq:c1-1-5} and part (a) of
Theorem~\ref{thm:c1-1-1} to
\[
|\hat{g}(x)-\hat{f}(x)(e^{2\pi i(h\cdot x)}-1)/h_{k}|,
\]
and then letting $h_{k} \to 0$ we obtain:

\begin{theorem}\label{thm:c1-1-8}
If $f \in L^{1}(E_{n})$ and $g$ is the partial derivative of $f$ with respect to
$x_{k}$ in the $L^{1}$ norm then
\[
\hat{g}(x)=2\pi i x_{k}\hat{f}(x).
\]
\end{theorem}

Both Theorems~\ref{thm:c1-1-7} and \ref{thm:c1-1-8} can be extended to higher
derivatives. Without going into details, we note the following formulas:
\begin{equation}\label{eq:c1-1-9}
\begin{aligned}
\textup{(i)}\quad &P(D)\hat{f}(x)=(P(-2\pi i t)f(t))^{\wedge}(x),\\
\textup{(ii)}\quad &(P(D)f)^{\wedge}(x)=P(2\pi i x)\hat{f}(x),
\end{aligned}
\end{equation}
where, for an $n$-tuple $\alpha = (\alpha_{1}, \ldots, \alpha_{n})$ of
nonnegative integers we let
$x^{\alpha} = x_{1}^{\alpha_{1}} x_{2}^{\alpha_{2}} \cdots x_{n}^{\alpha_{n}}$,
$D^{\alpha} = \partial^{\alpha_{1} + \alpha_{2} + \cdots + \alpha_{n}} / \partial x_{1}^{\alpha_{1}} \partial x_{2}^{\alpha_{2}} \cdots \partial x_{n}^{\alpha_{n}}$,
$P$ is a polynomial in the $n$ variables $x_{1}$, $x_{2}$, $\ldots$, $x_{n}$
and $P(D)$ is the associated differential operator (i.e., we replace $x^{\alpha}$
by $D^{\alpha}$ in $P(x)$).

We now turn to the problem of inverting the Fourier transform. That is, we
shall consider the question: Given the Fourier transform $\hat{f}$ of an
integrable function $f$, how do we obtain $f$ back again from $\hat{f}$? The
reader familiar with the elementary theory of Fourier series and integrals would
expect $f(t)$ to be equal to the integral
\begin{equation}\label{eq:c1-1-10}
\int_{E_{n}}\hat{f}(x)e^{2\pi i t\cdot x}\,dx.
\end{equation}

Unfortunately, $\hat{f}$ need not be integrable (for example, let $n = 1$ and
$f$ be the characteristic function of a finite interval). In order to get around
this difficulty we shall use certain summability methods for integrals. We first
introduce the Abel method of summability\index{Abel method of summability@Abel method of summability}, whose analog for series is very
well-known. For each $\varepsilon > 0$ we define the $Abel$ $mean$
$A_{\varepsilon} = A_{\varepsilon}(f)$ to be the integral
\[
A_{\varepsilon}(f)=A_{\varepsilon}=\int_{E_{n}}f(x)e^{-\varepsilon|x|}\,dx.
\]

It is clear that if $f \in L^{1}(E_{n})$ then
$\lim_{\varepsilon \to 0} A_{\varepsilon}(f) = \int_{E_{n}} f(x) \, dx$. On the
other hand, these Abel means are well-defined even when $f$ is not integrable
(if we only assume, for example, that $f$ is bounded, then
$A_{\varepsilon}(f)$ is defined for all $\varepsilon > 0$). Moreover, their
limit
\begin{equation}\label{eq:c1-1-11}
\lim_{\varepsilon\to0}A_{\varepsilon}(f)
=\lim_{\varepsilon\to0}\int_{E_{n}}f(x)e^{-\varepsilon|x|}\,dx
\end{equation}
may exist even when $f$ is not integrable. A classical example of such a case is
obtained by letting $f(x) = \sin x/x$ when $n = 1$.%
\footnote{As is well known, in this case
$\lim_{p\to\infty}\int_{0}^{p}f(x)\,dx$ exists. It is an easy exercise to show
that whenever $f$ is locally integrable and such a limit, $l$, exists the Abel
means $A_{\varepsilon}=\int_{0}^{\infty}e^{-\varepsilon x}f(x)\,dx$ converge to
$l$.}

Whenever the limit in \eqref{eq:c1-1-11} exists and is finite we say that
$\int_{E_{n}} f$ is \textit{Abel} summable to this limit.

A somewhat similar method of summability is Gauss summability\index{Gauss summability@Gauss summability}. This method is
defined by the Gauss (sometimes called Gauss-Weierstrass) means
\[
G_{\varepsilon}(f)=G_{\varepsilon}=\int_{E_{n}}f(x)e^{-\varepsilon|x|^{2}}\,dx.
\]
We say that $\int_{E_{n}} f$ is Gauss summable (to $l$) if
\begin{equation}\label{eq:c1-1-11p}
\tag{1.11$'$}
\lim_{\varepsilon\to0}G_{\varepsilon}(f)
=\lim_{\varepsilon\to0}\int_{E_{n}}f(x)e^{-\varepsilon|x|^{2}}\,dx
\end{equation}
exists and equals the number $l$.

We see that both \eqref{eq:c1-1-11} and \eqref{eq:c1-1-11p} can be put in the
form
\begin{equation}\label{eq:c1-1-12}
M_{\varepsilon,\Phi}(f)=M_{\varepsilon}(f)=\int_{E_{n}}\Phi(\varepsilon x)f(x)\,dx,
\end{equation}
where $\Phi \in C_{0}$ and $\Phi(0) = 1$. Then $\int_{E_{n}} f$ is summable to
$l$ if $\lim_{\varepsilon \to 0} M_{\varepsilon}(f) = l$. We shall call
$M_{\varepsilon}(f)$ the $\Phi$ means of this integral.

We shall need the Fourier transforms of the functions $e^{-\varepsilon|x|^{2}}$
and $e^{-\varepsilon|x|}$. The calculation of the first one is easy and is
reducible to the one dimensional case: Since
\begin{align*}
\int_{E_{n}}e^{-4\pi^{2}|x|^{2}}e^{-2\pi i x\cdot t}\,dx
&=\prod_{j=1}^{n}\int_{-\infty}^{\infty}e^{-4\pi^{2}x_{j}^{2}}e^{-2\pi i x_{j}t_{j}}\,dx_{j}\\
&=\prod_{j=1}^{n}\frac{1}{2\sqrt{\pi}}e^{-t_{j}^{2}/4}=2^{-n}\pi^{-n/2}e^{-|t|^{2}/4},
\end{align*}
the change of variables $(\sqrt{\alpha}/2\sqrt{\pi})y = x$ gives us

\begin{theorem}\label{thm:c1-1-13}
For all $\alpha > 0$ we have
\[
\int_{E_{n}}e^{-\pi\alpha|y|^{2}}e^{-2\pi i t\cdot y}\,dy
=\alpha^{-n/2}e^{-\pi|t|^{2}/\alpha}.%
\footnote{We note that this theorem shows that $e^{-\pi|x|^{2}}$ is its own
Fourier transform.}
\]
\end{theorem}

The second Fourier transform is somewhat harder to obtain:

\begin{theorem}\label{thm:c1-1-14}
For all $\alpha > 0$ we have
\[
\int_{E_{n}}e^{-2\pi|y|\alpha}e^{-2\pi i t\cdot y}\,dy
=c_{n}\frac{\alpha}{(\alpha^{2}+|t|^{2})^{(n+1)/2}},
\]
where $c_{n} = \Gamma[(n + 1)/2]/(\pi^{(n+1)/2})$\index{cn (normalization of Poisson kernel)@$c_n$ (normalization of Poisson kernel)}.
\end{theorem}

\begin{proof}
By a change of variables we see that it suffices to show this result when
$\alpha = 1$. In order to show this let us momentarily assume that
\[
\textup{(i)}\qquad
e^{-\beta}=\frac{1}{\sqrt{\pi}}\int_{0}^{\infty}\frac{e^{-u}}{\sqrt{u}}e^{-\beta^{2}/4u}\,du
\]
whenever $\beta > 0$. Then, using Theorem~\ref{thm:c1-1-13} to establish the
third equality,
\begin{align*}
\int_{E_{n}}e^{-2\pi|y|}e^{-2\pi i t\cdot y}\,dy
&=\int_{E_{n}}\left\{\frac{1}{\sqrt{\pi}}\int_{0}^{\infty}\frac{e^{-u}}{\sqrt{u}}e^{-4\pi^{2}|y|^{2}/4u}\,du\right\}e^{-2\pi i t\cdot y}\,dy\\
&=\frac{1}{\sqrt{\pi}}\int_{0}^{\infty}\frac{e^{-u}}{\sqrt{u}}\left\{\int_{E_{n}}e^{-4\pi^{2}|y|^{2}/4u}e^{-2\pi i t\cdot y}\,dy\right\}\,du\\
&=\frac{1}{\sqrt{\pi}}\int_{0}^{\infty}\frac{e^{-u}}{\sqrt{u}}\left\{\left(\sqrt{\frac{u}{\pi}}\right)^{n}e^{-u|t|^{2}}\right\}\,du\\
&=\frac{1}{\pi^{(n+1)/2}}\int_{0}^{\infty}e^{-u}u^{(n-1)/2}e^{-u|t|^{2}}\,du\\
&=\frac{1}{\pi^{(n+1)/2}}\frac{1}{(1+|t|^{2})^{(n+1)/2}}\int_{0}^{\infty}e^{-s}s^{(n-1)/2}\,ds\\
&=\frac{\Gamma[(n+1)/2]}{\pi^{(n+1)/2}}\frac{1}{(1+|t|^{2})^{(n+1)/2}}.
\end{align*}

Consequently the theorem will be established once we show (i). We do this by
using the identities
\[
\textup{(ii)}\qquad
e^{-\beta}=\frac{2}{\pi}\int_{0}^{\infty}\frac{\cos\beta x}{1+x^{2}}\,dx,\qquad\beta>0
\qquad\text{and}
\]
\[
\textup{(iii)}\qquad
\frac{1}{1+x^{2}}=\int_{0}^{\infty}e^{-(1+x^{2})u}\,du.
\]
The second of these equalities is obvious, while the first is an easy
application of the theory of residues to the function $e^{i\beta z}/(1 + z^{2})$.
Thus,
\begin{align*}
e^{-\beta}&=\frac{2}{\pi}\int_{0}^{\infty}\frac{\cos\beta x}{1+x^{2}}\,dx\\
&=\frac{2}{\pi}\int_{0}^{\infty}\cos\beta x\Biggl\{\int_{0}^{\infty}e^{-u}e^{-ux^{2}}\,du\Biggr\}dx\\
&=\frac{2}{\pi}\int_{0}^{\infty}e^{-u}\Biggl\{\int_{0}^{\infty}e^{-ux^{2}}\cos\beta x\,dx\Biggr\}du\\
&=\frac{2}{\pi}\int_{0}^{\infty}e^{-u}\Biggl\{\frac{1}{2}\int_{-\infty}^{\infty}e^{-ux^{2}}e^{i\beta x}\,dx\Biggr\}du\\
&=\frac{2}{\pi}\int_{0}^{\infty}e^{-u}\Biggl\{\pi\int_{-\infty}^{\infty}e^{-4\pi^{2}uy^{2}}e^{-2\pi i\beta y}\,dy\Biggr\}du\\
&=\frac{2}{\pi}\int_{0}^{\infty}e^{-u}\Biggl\{\frac{1}{2}\sqrt{\frac{\pi}{u}}e^{-\beta^{2}/4u}\Biggr\}du.
\end{align*}
and the theorem is proved.
\end{proof}

We shall denote the Fourier transforms of $e^{-4\pi^{2}\alpha|y|^{2}}$ and
$e^{-2\pi\alpha|y|}$, $\alpha > 0$, by $W$ and $P$, respectively. That is,
$W(t, \alpha) = (4\pi\alpha)^{-(n/2)}e^{-|t|^{2}/4\alpha}$ and
$P(t, \alpha) = c_{n}[\alpha/(\alpha^{2} + |t|^{2})^{(n+1)/2}]$. The first of
these two functions is called the Weierstrass (or Gauss--Weierstrass) kernel
while the second is called the Poisson kernel\index{Poisson kernel@Poisson kernel}.%
\footnote{In later sections we shall encounter other Poisson kernels; these
will be associated with certain domains. For this reason we shall, later on,
call the kernel we just introduced the Poisson kernel associated with the upper
half-space
$E_{n+1}^{+}=\{x=(x_{1},\ldots,x_{n},x_{n+1})\in E_{n+1};x_{n+1}>0\}$.}

We shall show, in particular, that the Abel and Gauss means of the integral
\eqref{eq:c1-1-10} converge to $f$ (both in the norm and a.e.). This will be
done by deriving a general formula from which we will be able to express the
Abel and Gauss means of this integral in terms of convolutions of $f$ with the
Poisson and Weierstrass kernels\index{Weierstrass kernel@Weierstrass kernel}. In order to obtain these expressions we will
use the following important result.

\begin{theorem}[The Multiplication Formula]\label{thm:c1-1-15}\index{multiplication formula@multiplication formula}
If $f$ and $g$ belong to $L^{1}(E_{n})$ then
\[
\int_{E_{n}}\hat{f}(x)g(x)\,dx=\int_{E_{n}}f(x)\hat{g}(x)\,dx.
\]
\end{theorem}

\begin{proof}
Applying Fubini's theorem, we have
\begin{align*}
\int_{E_{n}}\hat{f}(x)g(x)\,dx
&=\int_{E_{n}}\left\{\int_{E_{n}}f(t)e^{-2\pi i t\cdot x}\,dt\right\}g(x)\,dx\\
&=\int_{E_{n}}\left\{\int_{E_{n}}g(x)e^{-2\pi i t\cdot x}\,dx\right\}f(t)\,dt\\
&=\int_{E_{n}}f(t)\hat{g}(t)\,dt.
\end{align*}
\end{proof}

Suppose now that the function $\Phi$ in \eqref{eq:c1-1-12} is integrable and its
Fourier transform is $\hat{\Phi} = \varphi$. If we let
$\varphi_{\varepsilon}(x) = \varepsilon^{-n}\varphi(x/\varepsilon)$ for
$\varepsilon > 0$, then \eqref{eq:c1-1-6} implies
$(\delta_{\varepsilon}\Phi)^{\wedge}(x) = \varepsilon^{-n}\varphi(x/\varepsilon) = \varphi_{\varepsilon}(x)$.
For example, Theorem~\ref{thm:c1-1-13} shows us that when
$\Phi(x) = e^{-4\pi^{2}|x|^{2}}$ we have
$\varphi_{\varepsilon}(x) = W(x, \varepsilon^{2})$ and, similarly from
Theorem~\ref{thm:c1-1-14} we see that $\varphi_{\varepsilon}(x) = P(x, \varepsilon)$
when $\Phi(x) = e^{-2\pi|x|}$. Then, applying the multiplication formula
(Theorem~\ref{thm:c1-1-15}) to $f(x)$ and
$e^{2\pi i t \cdot x} \delta_{\varepsilon}\Phi(x)$ and using \eqref{eq:c1-1-5},
part (ii), in order to evaluate the Fourier transform of the latter function we
obtain

\begin{theorem}\label{thm:c1-1-16}
If $f$ and $\Phi$ belong to $L^{1}(E_{n})$ and $\varphi = \hat{\Phi}$ then
\[
\int_{E_{n}}\hat{f}(x)e^{2\pi i t\cdot x}\Phi(\varepsilon x)\,dx
=\int_{E_{n}}f(x)\varphi_{\varepsilon}(x-t)\,dx
\]
for all $\varepsilon > 0$. In particular,
\[
\int_{E_{n}}\hat{f}(x)e^{2\pi i t\cdot x}e^{-2\pi\varepsilon|x|}\,dx
=\int_{E_{n}}f(x)P(x-t,\varepsilon)\,dx
\]
and
\[
\int_{E_{n}}\hat{f}(x)e^{2\pi i t\cdot x}e^{-4\pi^{2}\alpha|x|^{2}}\,dx
=\int_{E_{n}}f(x)W(x-t,\alpha)\,dx
\]
for all $\alpha > 0$.
\end{theorem}

We shall now show that the integral \eqref{eq:c1-1-10} is summable to $f$ for a
large class of methods of summability that includes both Abel and Gauss
summability. We shall first show that the means
$\int_{E_{n}} \hat{f}(x) e^{2\pi i t \cdot x} \Phi(\varepsilon x) \, dx$
(see \eqref{eq:c1-1-12}) converge to $f$ in the $L^{1}$ norm under the following
very general conditions: $\Phi$ and $\hat{\Phi} = \varphi$ are both integrable
and $\int_{E_{n}} \varphi(x) \, dx = 1$. The following lemma will show that the
Poisson and Gauss--Weierstrass kernels satisfy these conditions.

\begin{lemma}\label{lem:c1-1-17}
\textup{(a)} $\int_{E_{n}} W(x, \alpha) \, dx = 1$ for all $\alpha > 0$;

\textup{(b)} $\int_{E_{n}} P(x, \varepsilon) \, dx = 1$ for all
$\varepsilon > 0$.
\end{lemma}

\begin{proof}
By a change of variable we first note that
$\int_{E_{n}} W(x, \alpha) \, dx = \int_{E_{n}} W(x, 1) \, dx$ and
$\int_{E_{n}} P(x, \varepsilon) \, dx = \int_{E_{n}} P(x, 1) \, dx$. Thus, it
suffices to prove the lemma when $\alpha = 1 = \varepsilon$. In this case, part
(a) is an immediate consequence of the one dimensional result
$\int_{-\infty}^{\infty} e^{-x^{2}/4} \, dx = 2\sqrt{\pi}$ (for the
$n$-dimensional integral can be written as the product of $n$ such integrals).
In order to establish (b) we first note that
$1/c_{n} = \pi^{(n+1)/2}/\Gamma[(n+1)/2]$ is half the surface area of the unit
sphere $\Sigma_{n}$ of $E_{n+1}$. Denoting this surface area by $\omega_{n}$, we
see that (b) is then equivalent to
\[
\int_{E_{n}}\frac{dx}{(1+|x|^{2})^{(n+1)/2}}=\frac{\omega_{n}}{2}.
\]
But, letting $r = |x|$, $x' = x/r$ (when $x \neq 0$),
$\Sigma_{n-1} = \{x \in E_{n}; |x| = 1\}$, $dx'$ the element of surface area on
$\Sigma_{n-1}$ and, finally, putting $r = \tan \theta$, we have
\begin{align*}
\int_{E_{n}}\frac{dx}{(1+|x|^{2})^{(n+1)/2}}
&=\int_{0}^{\infty}\left(\int_{\Sigma_{n-1}}\frac{1}{(1+r^{2})^{(n+1)/2}}\,dx'\right)r^{n-1}\,dr\\
&=\omega_{n-1}\int_{0}^{\infty}\frac{r^{n-1}}{(1+r^{2})^{(n+1)/2}}\,dr\\
&=\omega_{n-1}\int_{0}^{\pi/2}\sin^{n-1}\theta\,d\theta.%
\footnote{We are, of course, expressing our original integral in
``polar coordinates'' $(r, x')$. We remind the reader that
$dx = r^{n-1} dr dx'$ (see \eqref{eq:c1-4-14} and the bibliographical notes for
further details).}
\end{align*}
But $\omega_{n-1}\sin^{n-1}\theta$ is clearly the surface area of the sphere of
radius $\sin\theta$ obtained by intersecting $\Sigma_{n}$ with the hyperplane
$x_{n} = \cos\theta$. Thus, the area of the upper half of $\Sigma_{n}$ is
obtained by summing these $(n-1)$ dimensional areas as $\theta$ ranges from $0$
to $\pi/2$; that is,
\[
\omega_{n-1}\int_{0}^{\pi/2}\sin^{n-1}\theta\,d\theta=\frac{\omega_{n}}{2},
\]
which is the desired result.
\end{proof}

\begin{theorem}\label{thm:c1-1-18}
Suppose $\varphi \in L^{1}(E_{n})$ with
$\int_{E_{n}} \varphi(x) \, dx = 1$ and for $\varepsilon > 0$ let
$\varphi_{\varepsilon}(x) = \varepsilon^{-n} \varphi(x/\varepsilon)$. If
$f \in L^{p}(E_{n})$, $1 \leq p < \infty$, or
$f \in C_{0} \subset L^{\infty}(E_{n})$, then
$\|f * \varphi_{\varepsilon} - f\|_{p} \to 0$ as $\varepsilon \to 0$. In
particular, $u(x, \varepsilon) = \int_{E_{n}} f(t) P(x - t, \varepsilon) \, dt$
and $s(x, \varepsilon) = \int_{E_{n}} f(t) W(x - t, \varepsilon) \, dt$ converge
to $f$ in the $L^{p}$ norm as $\varepsilon \to 0$.%
\footnote{The function $u(x, \varepsilon)$, defined for $x \in E_{n}$ and
$\varepsilon > 0$, is called the Poisson integral of $f$; $s(x, \varepsilon)$ is
known as the Gauss--Weierstrass integral of $f$.}
\end{theorem}

\begin{proof}
We have, by a change of variables,
\[
\int_{E_{n}}\varphi_{\varepsilon}(t)\,dt
=\int_{E_{n}}\varepsilon^{-n}\varphi(t/\varepsilon)\,dt
=\int_{E_{n}}\varphi(t)\,dt=1.
\]
Hence,
\begin{align*}
(f*\varphi_{\varepsilon})(x)-f(x)
&=\int_{E_{n}}f(x-t)\varphi_{\varepsilon}(t)\,dt-\int_{E_{n}}f(x)\varphi_{\varepsilon}(t)\,dt\\
&=\int_{E_{n}}[f(x-t)-f(x)]\varphi_{\varepsilon}(t)\,dt.
\end{align*}
Therefore, by Minkowski's inequality for integrals\index{Poisson integral@Poisson integral}\index{Weierstrass integral@Weierstrass integral}
\begin{align*}
\|f*\varphi_{\varepsilon}-f\|_{p}
&\leq\int_{E_{n}}\left(\int_{E_{n}}|f(x-t)-f(x)|^{p}\,dx\right)^{1/p}\varepsilon^{-n}|\varphi(t/\varepsilon)|\,dt\\
&=\int_{E_{n}}\left(\int_{E_{n}}|f(x-\varepsilon t)-f(x)|^{p}\,dx\right)^{1/p}|\varphi(t)|\,dt.
\end{align*}
The expression
$(\int_{E_{n}}|f(x+h)-f(x)|^{p}\,dx)^{1/p}=\omega_{p,f}(h)=\omega(h)$ is known
as the $L^{p}$ modulus of continuity\index{modulus of continuity@modulus of continuity!Lp@$L^{p}$} of $f$. It is obviously bounded as a
function of $h$ since $\omega(h)\leq2\|f\|_{p}$. Moreover, $\omega(h)\to0$ as
$|h|\to0$. This is clearly the case when $f$ is continuous and has compact
support. The general case is easily obtained by approximating $f$ (in the
$L^{p}$ norm) by such functions.

We have shown, therefore, that
\[
\|f*\varphi_{\varepsilon}-f\|_{p}\leq\int_{E_{n}}\omega(\varepsilon t)|\varphi(t)|\,dt.
\]
But, by the Lebesgue dominated convergence theorem, the integral on the right
tends to $0$ with $\varepsilon$ since the integrand tends to $0$ and is
majorized by the integrable function $2\|f\|_{p}|\varphi(t)|$. This proves the
theorem.
\end{proof}

Though it is not related to our study of the inversion of the Fourier integral
we observe that essentially the same argument gives us the following result,
which will be useful to us in the future:

\begin{corollary}\label{cor:c1-1-19}
Suppose $\varphi \in L^{1}(E_{n})$ and $\int_{E_{n}} \varphi(t) \, dt = 0$ then
$\|f * \varphi_{\varepsilon}\|_{p} \to 0$ as $\varepsilon \to 0$ whenever
$f \in L^{p}(E_{n})$, $1 \leq p < \infty$, or
$f \in C_{0} \subset L^{\infty}(E_{n})$.
\end{corollary}

\begin{proof}
Once we observe that
$(f*\varphi_{\varepsilon})(x)=(f*\varphi_{\varepsilon})(x)-f(x)\cdot0=(f*\varphi_{\varepsilon})(x)-f(x)\int_{E_{n}}\varphi_{\varepsilon}(t)\,dt=\int_{E_{n}}[f(x-t)-f(x)]\varphi_{\varepsilon}(t)\,dt$,
the rest of the argument is precisely that used in the last proof.
\end{proof}

From Theorems~\ref{thm:c1-1-16} and \ref{thm:c1-1-18} we obtain the following
solution to the Fourier inversion problem:

\begin{theorem}\label{thm:c1-1-20}
If $\Phi$ and its Fourier transform $\varphi = \hat{\Phi}$ are integrable and
$\int_{E_{n}} \varphi(x) \, dx = 1$, then the $\Phi$ means of the integral
$\int_{E_{n}} \hat{f}(t) e^{2\pi i t \cdot x} \, dt$ converge to $f(x)$ in the
$L^{1}$ norm. In particular, the Abel and Gauss means of this integral converge
to $f(x)$ in the $L^{1}$ norm.
\end{theorem}

Since
$s(x, \alpha) = \int_{E_{n}} \hat{f}(t) e^{2\pi i x \cdot t} e^{-4\pi^{2} \alpha |t|^{2}} \, dt$
converges in $L^{1}$ to $f(x)$ as $\alpha > 0$ tends to $0$, we can find a
sequence $\alpha_{k} \to 0$ such that $s(x, \alpha_{k}) \to f(x)$ for a.e. $x$.
If we further assume that $\hat{f} \in L^{1}(E_{n})$ the Lebesgue dominated
convergence theorem gives us the following pointwise equality:

\begin{corollary}\label{cor:c1-1-21}
If both $f$ and $\hat{f}$ are integrable then
\[
f(x)=\int_{E_{n}}\hat{f}(t)e^{2\pi i x\cdot t}\,dt
\]
for almost every $x$.%
\footnote{We know from Theorem~\ref{thm:c1-1-1} that $\hat{f}$ is continuous.
If $\hat{f}$ is integrable, the integral
$\int_{E_{n}}\hat{f}(t)e^{2\pi i x\cdot t}\,dt$ also defines a continuous
function (in fact, it equals $(\hat{f})^{\wedge}(-x)$). Thus, by changing $f$ on
a set of measure $0$ we can obtain equality in Corollary~\ref{cor:c1-1-21} for
all $x$.}
\end{corollary}

We have singled out the Gauss-Weierstrass and the Abel methods of summability.
The former is probably the simplest; we shall see in the second chapter that the
latter is intimately connected with harmonic functions and provides us with very
powerful tools in Fourier analysis.%
\footnote{In the same way, the Gauss--Weierstrass method is intimately connected
with the solution of the heat equation.}

A typical example of the relative simplicity of the Gauss--Weierstrass method is
Theorem~\ref{thm:c1-1-13} whose proof is much simpler than the corresponding
result for the Abel method (Theorem~\ref{thm:c1-1-14}).

We also have seen that part (a) of Lemma~\ref{lem:c1-1-17} was much easier to
prove than part (b). It is of interest to remark that
Corollary~\ref{cor:c1-1-21} enables us to derive the latter without making the
calculation occurring in the proof of Lemma~\ref{lem:c1-1-17}. In fact, applying
Corollary~\ref{cor:c1-1-21} to $f(x) = e^{-2\pi\varepsilon|x|}$ and using
Theorem~\ref{thm:c1-1-14} we obtain
$\int_{E_{n}} P(t, \varepsilon)e^{2\pi i x \cdot t} \, dt = e^{-2\pi\varepsilon|x|}$.
Putting $x = 0$ we obtain $\int_{E_{n}} P(t, \varepsilon) \, dt = 1$.

It is clear from Theorem~\ref{thm:c1-1-20} that if $\hat{f}(x) = 0$ for all $x$
then $f(t) = 0$ for almost every $t$. Applying this to $f = f_{1} - f_{2}$ we
obtain the following uniqueness result for the Fourier transform:

\begin{corollary}\label{cor:c1-1-22}
If $f_{1}$ and $f_{2}$ belong to $L^{1}(E_{n})$ and
$\hat{f}_{1}(x) = \hat{f}_{2}(x)$ for $x \in E_{n}$ then $f_{1}(t) = f_{2}(t)$
for almost every $t \in E_{n}$.
\end{corollary}

The Fourier inversion problem also has a solution in the pointwise sense. We
shall show, in particular, that the Abel and Gauss means of the integral\index{differentiability of the integral@differentiability of the integral}
\eqref{eq:c1-1-10} converge almost everywhere to $f(t)$. In order to do this we
need to introduce some new concepts and notation. We first state an important
result in the theory of differentiation of integrals of functions defined on
$E_{n}$. If $f$ is locally integrable on $E_{n}$ then for almost every
$x \in E_{n}$
\begin{equation}\label{eq:c1-1-23}
\frac{1}{r^{n}}\int_{|t|<r}[f(x-t)-f(x)]\,dt\to0
\end{equation}
as $r \to 0$.%
\footnote{If we let $S(x;r)$ denote the sphere
$\{t\in E_{n};|t-x|<r\}$ and $|S(x;r)|$ its measure, then \eqref{eq:c1-1-23} is
equivalent to the equality
$f(x)=\lim_{r\to0}|S(x;r)|^{-1}\int_{S(x;r)}f(t)\,dt$. When $n=1$ this becomes
$f(x)=\lim_{r\to0}(1/2r)\int_{x-r}^{x+r}f(t)\,dt$ and it is a classical theorem
of Lebesgue that this equality holds a.e. We defer the proof of the fact that
\eqref{eq:c1-1-23} holds a.e. to the second chapter where we shall see that it is
a consequence of a general technique involving the Hardy-Littlewood maximal
function (see Corollary \ref{cor:c2-3-14} in that chapter).}
In particular, this is true for $f \in L^{p}(E_{n})$. We shall call the set of
points $x$ for which \eqref{eq:c1-1-23} holds the set of points where the
integral of $f$ is differentiable. This set is naturally associated with each
locally integrable function $f$. There exists, a smaller but closely related set
that (as we shall see) is also naturally associated with $f$. This is the set of
all $x$ such that
\begin{equation}\label{eq:c1-1-24}
\frac{1}{r^{n}}\int_{|t|<r}|f(x-t)-f(x)|\,dt\to0
\end{equation}
as $r \to 0$. This set is called the Lebesgue set\index{Lebesgue set@Lebesgue set} of $f$ and it includes almost
all points of $E_{n}$. To see this we note that since the function
$|f(t) - \rho|$ is locally integrable it follows from \eqref{eq:c1-1-23} that the
set $F_{\rho}$ of all $x \in E_{n}$ for which it is not true that
\[
\lim_{r\to0}\frac{1}{r^{n}}\int_{|t|<r}\{|f(x-t)-\rho|-|f(x)-\rho|\}\,dt=0
\]
has measure $0$. Thus, the union, $F$, over all rational numbers $\rho$, of the
sets $F_{\rho}$ has measure $0$. We claim that if $x \in E_{n} - F$ then
\eqref{eq:c1-1-24} holds. In fact, given such an $x$ and an $\varepsilon > 0$
let $\rho$ be a rational number such that $|f(x) - \rho| < \varepsilon/2$ and let
$\Omega_{n}$ denote the volume of the (solid) unit sphere in $E_{n}$. Then,
\begin{align*}
\frac{1}{\Omega_{n} r^{n}}\int_{|t|<r}|f(x-t)-f(x)|\,dt
&\leq\frac{1}{\Omega_{n} r^{n}}\int_{|t|<r}|f(x-t)-\rho|\,dt\\
&\quad+\frac{1}{\Omega_{n} r^{n}}\int_{|t|<r}|\rho-f(x)|\,dt.
\end{align*}
The first term on the right tends to $|f(x) - \rho| < \varepsilon/2$ as
$r \to 0$, while the second term equals $|\rho - f(x)| < \varepsilon/2$. Thus,
for $r$ close to $0$,
\[
(\Omega_{n}r^{n})^{-1}\int_{|t|<r}|f(x-t)-f(x)|\,dt<\varepsilon
\]
and it follows that almost all points $x$ of $E_{n}$ belong to the Lebesgue set
of $f$.

We shall show that under suitable conditions the integral \eqref{eq:c1-1-10} is
$\Phi$-summable to $f(t)$ whenever $t$ belongs to the Lebesgue set of $f$. More
precisely, the following general result, together with
Theorem~\ref{thm:c1-1-16}, gives us conditions on the Fourier transform,
$\varphi$, of $\Phi$ for this to be the case:

\begin{theorem}\label{thm:c1-1-25}
Suppose $\varphi \in L^{1}(E_{n})$. Let
$\psi(x) = \operatorname{ess.\ sup.}_{|t| \geq |x|}|\varphi(t)|$ and, for
$\varepsilon > 0$, let $\varphi_{\varepsilon}(x) = \varepsilon^{-n}\varphi(x/\varepsilon)$.
If $\psi \in L^{1}(E_{n})$ and $f \in L^{p}(E_{n})$, $1 \leq p \leq \infty$, then
$\lim_{\varepsilon \to 0}(f * \varphi_{\varepsilon})(x) = f(x) \int_{E_{n}} \varphi(t) \, dt$
whenever $x$ belongs to the Lebesgue set of $f$. In particular, the Poisson and
Gauss--Weierstrass integrals of $f$,
$\int_{E_{n}} f(t)P(x - t, \varepsilon) \, dt$ and
$\int_{E_{n}} f(t)W(x - t, \varepsilon) \, dt$, converge to $f(x)$ as
$\varepsilon \to 0$ for almost every $x \in E_{n}$.
\end{theorem}

\begin{proof}
We fix a point $x$ of the Lebesgue set of $f$ and choose $\delta > 0$. Then we
can find an $\eta > 0$ such that
\[
\textup{(i)}\qquad
r^{-n}\int_{|t|<r}|f(x-t)-f(x)|\,dt<\delta
\]
provided $r \leq \eta$.

We have, since $\int_{E_{n}}\varphi_{\varepsilon}(t)\,dt=\int_{E_{n}}\varphi(t)\,dt=a$
for all $\varepsilon>0$,
\begin{align*}
|(f*\varphi_{\varepsilon})(x)-af(x)|
&=\left|\int_{E_{n}}\{f(x-t)-f(x)\}\varphi_{\varepsilon}(t)\,dt\right|\\
&\leq\left|\int_{|t|<\eta}\{f(x-t)-f(x)\}\varphi_{\varepsilon}(t)\,dt\right|\\
&\quad+\left|\int_{|t|\geq\eta}\{f(x-t)-f(x)\}\varphi_{\varepsilon}(t)\,dt\right|\\
&=I_{1}+I_{2}.
\end{align*}

In order to estimate $I_{1}$ we first observe that the function $\psi$ is radial\index{radial function@radial function}
(i.e., $\psi(x_{1})=\psi(x_{2})$ if $|x_{1}|=|x_{2}|$) and, if we put
$\psi_{0}(r)=\psi(x)$ when $|x|=r$, then $\psi_{0}$ is a decreasing function of
$r$. Thus, if $\Omega_{n}$ denotes the volume of the solid unit sphere of
$E_{n}$, we have
\[
\frac{\Omega_{n}(2^{n}-1)}{2^{n}}r^{n}\psi_{0}(r)\leq\int_{r/2\leq|x|\leq r}\psi(x)\,dx\to0
\]
as $r \to 0$ or as $r \to \infty$. That is, $\lim r^{n} \psi_{0}(r) = 0$ as $r$
tends to $0$ or $\infty$. In particular, there exists a constant $A$ such that
$r^{n} \psi_{0}(r) \leq A$ for $0 < r < \infty$. If $\Sigma_{n-1}$ denotes the
surface of the unit sphere $\{t' \in E_{n}; |t'| = 1\}$, we let
$g(r) = \int_{\Sigma_{n-1}} |f(x - rt') - f(x)| \, dt'$, where $dt'$ is the
element of surface area on $\Sigma_{n-1}$. Condition (i) is therefore equivalent
to
\[
G(r)=\int_{0}^{r}s^{n-1}g(s)\,ds\leq\delta r^{n}
\]
provided $r \leq \eta$. Using this notation and these observations, therefore,
we have
\begin{align*}
I_{1}&\leq\int_{|t|<\eta}|f(x-t)-f(x)|\varepsilon^{-n}\psi(t/\varepsilon)\,dt
=\int_{0}^{\eta}r^{n-1}g(r)\varepsilon^{-n}\psi_{0}(r/\varepsilon)\,dr\\
&=G(r)\varepsilon^{-n}\psi_{0}(r/\varepsilon)\big]_{0}^{\eta}
-\int_{0}^{\eta}G(r)\,d(\varepsilon^{-n}\psi_{0}(r/\varepsilon))\\
&\leq\delta r^{n}\varepsilon^{-n}\psi_{0}(r/\varepsilon)\big]_{0}^{\eta}
-\int_{0}^{\eta/\varepsilon}G(\varepsilon s)\varepsilon^{-n}\,d\psi_{0}(s)\\
&\leq\delta A-\int_{0}^{\eta/\varepsilon}\delta s^{n}\,d\psi_{0}(s)
\leq\delta\Biggl(A-\int_{0}^{\infty}s^{n}\,d\psi_{0}(s)\Biggr).
\end{align*}
But
$-\int_{0}^{\infty}s^{n}\,d\psi_{0}(s)=n\int_{0}^{\infty}s^{n-1}\psi_{0}(s)\,ds=(n/\omega_{n-1})\int_{E_{n}}\psi(x)\,dx$.
This shows that there exists a constant $B$, depending only on $\psi$, such that
$I_{1}\leq B\delta$.

In order to estimate $I_{2}$ we let
$\psi_{\varepsilon}(x) = \varepsilon^{-n}\psi(x/\varepsilon)$ and we denote the
characteristic function of the set of all $x \in E_{n}$ such that $|x| \geq \eta$
by $\chi_{\eta}$. Then, if $(1/p) + (1/p') = 1$,
\[
I_{2}\leq\|f\|_{p}\|\chi_{\eta}\psi_{\varepsilon}\|_{p'}+|f(x)|\cdot\|\chi_{\eta}\psi_{\varepsilon}\|_{1}.
\]
Since
\[
\|\chi_{\eta}\psi_{\varepsilon}\|_{1}
=\int_{|x|\geq\eta}\psi_{\varepsilon}(x)\,dx
=\int_{|x|\geq\eta/\varepsilon}\psi(x)\,dx,
\]
the second summand tends to $0$ with $\varepsilon$. We can show that the same is
true of the first summand: Since $p' = 1 + (p'/p)$, an application of H\"older's
inequality gives us
\begin{align*}
\|\chi_{\eta}\psi_{\varepsilon}\|_{p'}
&=\left(\int_{|x|\geq\eta}[\psi_{\varepsilon}(x)]^{p'}\,dx\right)^{1/p'}
=\left(\int_{|x|\geq\eta}\psi_{\varepsilon}(x)[\psi_{\varepsilon}(x)]^{p'/p}\,dx\right)^{1/p'}\\
&\leq\|\chi_{\eta}\psi_{\varepsilon}\|_{\infty}^{1/p}\|\chi_{\eta}\psi_{\varepsilon}\|_{1}^{1/p'}.
\end{align*}
But, as we have already noted,
$\|\chi_{\eta}\psi_{\varepsilon}\|_{\infty} = \sup_{|x| \geq \eta} \psi_{\varepsilon}(x) = \eta^{-n}(\eta/\varepsilon)^{n}\allowbreak\psi_{0}(\eta/\varepsilon) \to 0$
as $\varepsilon \to 0$.

We have shown, therefore, that for $\varepsilon$ sufficiently small
\[
|(f*\varphi_{\varepsilon})(x)-af(x)|
\]
is bounded by a constant multiple of $\delta$, this constant depending only on
$\psi$. This clearly proves the theorem.
\end{proof}

If $x$ is a point of continuity of $f$ then it certainly belongs to the Lebesgue
set of $f$. Consequently, Theorem~\ref{thm:c1-1-25} holds for any such point. If
$f$ is continuous at $0$, Theorems~\ref{thm:c1-1-25} and \ref{thm:c1-1-16} imply
\[
\lim_{\varepsilon\to0}\int_{E_{n}}\hat{f}(x)e^{-2\pi\varepsilon|x|}\,dx=f(0).
\]
If we further assume that $\hat{f} \geq 0$ it follows from Fatou's lemma that
$\hat{f} \in L^{1}(E_{n})$. Thus (see Theorem~\ref{thm:c1-1-1}, part (b)), the
integral \eqref{eq:c1-1-10} defines a continuous function of $t$ that equals
$f(t)$ almost everywhere by Corollary~\ref{cor:c1-1-21}. We therefore obtain the
following very useful result:

\begin{corollary}\label{cor:c1-1-26}
Suppose $f \in L^{1}(E_{n})$ and $\hat{f} \geq 0$. If $f$ is continuous at $0$
then $\hat{f}$ belongs to $L^{1}(E_{n})$ and
\[
f(t)=\int_{E_{n}}\hat{f}(x)e^{2\pi i t\cdot x}\,dx
\]
for almost every $x$. In particular,
\[
f(0)=\int_{E_{n}}\hat{f}(x)\,dx.
\]
\end{corollary}
% ============================================================
%  CHAPTER I -- The Fourier Transform   (printed pp. 16-35)
%  Part B: the tail of section 1 (Corollaries 1.27-1.28, printed p. 16),
%          sections 2, 3, 4 and the Bibliographical Notes.
%  Source : src/ch1.blocks.md, blocks from p28 to p47 (p48 is blank).
%  NOTE   : the first two corollaries below belong to section 1 and are
%           placed here only so that they directly follow Corollary 1.26
%           in ch01.tex; the shared counter therefore continues 1.27, 1.28
%           before \section{2} resets it.  Delete them if ch01.tex already
%           contains them.
% ============================================================

An immediate consequence of this corollary and the formulas in
Theorems~\ref{thm:c1-1-13} and \ref{thm:c1-1-14} is

\begin{corollary}\label{cor:c1-1-27}
\textup{(a)} $\int_{E_{n}} W(x, \alpha) e^{2\pi i t \cdot x} \, dx = e^{-4\pi^{2} \alpha |t|^{2}}$

\textup{(b)} $\int_{E_{n}} P(x, \alpha) e^{2\pi i t \cdot x} \, dx = e^{-2\pi\alpha|t|}$
for all $\alpha > 0$.
\end{corollary}

From Theorem~\ref{thm:c1-1-4} and Corollaries~\ref{cor:c1-1-22}
and \ref{cor:c1-1-27} we immediately obtain the semigroup properties\index{semigroup properties@semigroup properties} of the
Weierstrass and Poisson kernels.

\begin{corollary}\label{cor:c1-1-28}
If $\alpha_{1}$ and $\alpha_{2}$ are positive real numbers then
\[
\textup{(a)}\qquad W(x,\alpha_{1}+\alpha_{2})=\int_{E_{n}}W(x-t,\alpha_{1})W(t,\alpha_{2})\,dt
\]
and
\[
\textup{(b)}\qquad P(x,\alpha_{1}+\alpha_{2})=\int_{E_{n}}P(x-t,\alpha_{1})P(t,\alpha_{2})\,dt
\]
\end{corollary}

\section{The \texorpdfstring{$L^{2}$}{L2} Theory and the Plancherel Theorem}

The integral defining the Fourier transform is not defined in the Lebesgue
sense for the general function in $L^{2}(E_{n})$; nevertheless, the Fourier
transform has a natural definition on this space and a particularly elegant
theory.

If, in addition to being integrable, we assume $f$ to be square-integrable
then $\hat{f}$ will also be square-integrable. In fact, we have the following
basic result:

\begin{theorem}\label{thm:c1-2-1}
If $f \in L^{1} \cap L^{2}$ then $\|\hat{f}\|_{2} = \|f\|_{2}$.
\end{theorem}

\begin{proof}
Let $g(x) = \overline{f(-x)}$. Then, by Theorem~\ref{thm:c1-1-3},
$h = f * g \in L^{1}(E_{n})$ and, by Theorem~\ref{thm:c1-1-4},
$\hat{h} = \hat{f}\hat{g}$. But $\hat{g} = \overline{\hat{f}}$, thus,
$\hat{h} = |\hat{f}|^{2}$. Now, an application of
Corollary~\ref{cor:c1-1-26} shows that $\hat{h} \in L^{1}(E_{n})$ and
$h(0) = \int_{E_{n}} \hat{h}(x) \, dx$ (it is an immediate consequence of
Schwarz's inequality and the fact that the $L^{2}$ modulus of continuity
$\omega_{2}(f; \delta)$ tends to $0$ as $\delta \to 0$ that $h$, the
convolution of the two $L^{2}$ functions $f$ and $g$, is uniformly
continuous). We thus have
\begin{align*}
\int_{E_{n}}|\hat{f}|^{2}\,dx
&=\int_{E_{n}}\hat{h}\,dx=h(0)=\int_{E_{n}}f(x)g(0-x)\,dx\\
&=\int_{E_{n}}f(x)\overline{f(x)}\,dx=\int_{E_{n}}|f|^{2}\,dx.
\end{align*}
\end{proof}

This theorem asserts that the Fourier transform is a bounded linear operator
defined on the dense subset $L^{1} \cap L^{2}$ of $L^{2}(E_{n})$ (in fact, it
is an isometry). Therefore, there exists a unique bounded extension,
$\mathcal{F}$, of this operator to all of $L^{2}$. $\mathcal{F}$ will be
called the Fourier transform on $L^{2}$; we shall also use the notation
$\hat{f} = \mathcal{F}f$ whenever $f \in L^{2}(E_{n})$.

In general, if $f \in L^{2}(E_{n})$ this definition of the Fourier transform
gives us $\hat{f}$ as the $L^{2}$ limit of the sequence $\{\hat{h}_{k}\}$,
where $\{h_{k}\}$ is any sequence in $L^{1} \cap L^{2}$ converging to $f$ in
the $L^{2}$ norm. It is convenient to choose the sequence $\{h_{k}\}$, where
$h_{k}(t)$ equals $f(t)$ when $|t| \leq k$ and is zero elsewhere. Thus,
$\hat{f}$ is the $L^{2}$ limit of the sequence of functions $\hat{h}_{k}$
defined by
\begin{equation}\label{eq:c1-2-2}
\hat{h}_{k}(x)=\int_{|t|\leq k}f(t)e^{-2\pi i x\cdot t}\,dt
=\int_{E_{n}}h_{k}(t)e^{-2\pi i x\cdot t}\,dt.
\end{equation}

A linear operator on $L^{2}(E_{n})$ that is an isometry and maps onto
$L^{2}(E_{n})$ is called a unitary\index{unitary@unitary} operator. It is an immediate consequence of
Theorem~\ref{thm:c1-2-1} that $\mathcal{F}$ is an isometry. Moreover, we have
the additional property that $\mathcal{F}$ is onto:

\begin{theorem}\label{thm:c1-2-3}
The Fourier transform is a unitary operator on $L^{2}(E_{n})$.
\end{theorem}

\begin{proof}
Since $\mathcal{F}$ is an isometry its range is a closed subspace of
$L^{2}(E_{n})$. If this subspace were not all of $L^{2}(E_{n})$ we could find
a function $g$ such that $\int_{E_{n}} \hat{f}g \, dx = 0$ for all
$f \in L^{2}$ and $\|g\|_{2} \neq 0$. The multiplication formula
(Theorem~\ref{thm:c1-1-15}) obviously extends to $L^{2}$; consequently,
$\int_{E_{n}}\hat{f}g\,dx=\int_{E_{n}}f\hat{g}\,dx=0$ for all $f \in L^{2}$.
But this implies that $\hat{g}(x) = 0$ for almost every $x$, contradicting the
fact that $\|g\|_{2} = \|\hat{g}\|_{2} \neq 0$.
\end{proof}

Theorem~\ref{thm:c1-2-3} is a major part of the basic theorem in the $L^{2}$
theory of the Fourier transform:

\begin{theorem}\label{thm:c1-2-4}
The inverse of the Fourier transform, $\mathcal{F}^{-1}$, can be obtained by
letting $(\mathcal{F}^{-1}g)(x) = (\mathcal{F}g)(-x)$ for all
$g \in L^{2}(E_{n})$.
\end{theorem}

Theorems~\ref{thm:c1-2-3} and \ref{thm:c1-2-4} are usually referred to as the
Plancherel\index{Fourier transform on L2@Fourier transform on $L^2$} theorem\index{Plancherel theorem@Plancherel theorem}. Theorem~\ref{thm:c1-2-4} follows from the fact that
$\mathcal{F}^{-1}\hat{f}$ can be expressed as the $L^{2}$ limit of the
sequence defined by
\begin{equation}\label{eq:c1-2-5}
f_{k}(t)=\int_{|x|\leq k}\hat{f}(x)e^{2\pi i t\cdot x}\,dx.
\end{equation}

We first claim that this is clear when $\hat{f} \in L^{1} \cap L^{2}$ (a dense
subset of the range). To do this we only need to check that this expression
agrees with $\mathcal{F}^{*}\hat{f}$.%
\footnote{We consider $L^{2}(E_{n})$ as a Hilbert space with inner product
$(f, g) = \int_{E_{n}} f\bar{g} \, dx$ and assume the reader is familiar with
the fact that, if $T$ is unitary, $T^{-1}$ equals the adjoint\index{adjoint@adjoint} operator $T^{*}$
(i.e., the operator satisfying $(Tf, g) = (f, T^{*}g)$ for all $f$ and $g$ in
$L^{2}(E_{n})$). From this it follows that $T$ preserves the inner product:
$(Tf, Tg) = (f, g)$.}
But this is immediate, for letting
$\tilde{f}(t) = \int_{E_{n}} \hat{f}(x) e^{2\pi i t \cdot x} \, dx = \lim_{k \to \infty} f_{k}(t)$
(in $L^{2}$), we have
\begin{align*}
(g,\tilde{f})
&=\int_{E_{n}}g(t)\Biggl(\overline{\int_{E_{n}}\hat{f}(x)e^{2\pi i t\cdot x}\,dx}\Biggr)dt\\
&=\int_{E_{n}}\left(\int_{E_{n}}g(t)e^{-2\pi i x\cdot t}\,dt\right)\overline{\hat{f}(x)}\,dx
=(\mathcal{F}g,\hat{f})
\end{align*}
for all $g \in L^{1} \cap L^{2}$. That is,
$(g, \tilde{f}) = (\mathcal{F}g, \mathcal{F}f) = (g, f)$ for all
$g \in L^{1} \cap L^{2}$; therefore, $\tilde{f} = f$. The general case now
follows by the same considerations that were used to establish
\eqref{eq:c1-2-2}.

We see, therefore, that the problem of inverting the Fourier transform has a
very simple and elegant solution in $L^{2}$. In view of the material
developed in the last section it is natural to ask if there exists a solution
involving summability concepts as well. For example, since
$e^{-\varepsilon|x|}$, as a function of $x$, is square integrable, the Abel
means of the integral \eqref{eq:c1-1-10} are well defined. Is it true that
they converge to $f$ (in $L^{2}$ or a.e.)? The answer is ``yes'' and this
follows immediately from Theorems~\ref{thm:c1-1-18} and \ref{thm:c1-1-25} once
the identity in Theorem~\ref{thm:c1-1-16} is established for
$f \in L^{2}(E_{n})$. But the latter can be proved in the same way as
Theorem~\ref{thm:c1-1-16} by using the $L^{2}$ extension of the multiplication
formula (Theorem~\ref{thm:c1-1-15}).

Having defined the Fourier transform for functions in $L^{1}(E_{n})$ and
functions in $L^{2}(E_{n})$ we can easily define it on the class
$L^{1}(E_{n}) + L^{2}(E_{n})$ consisting of all functions $f = f_{1} + f_{2}$,
with $f_{1} \in L^{1}(E_{n})$ and $f_{2} \in L^{2}(E_{n})$. We do this by
putting $(f_{1} + f_{2})^{\wedge} = \hat{f}_{1} + \hat{f}_{2}$. If
$g_{1} + g_{2} = f_{1} + f_{2}$, with $g_{i} \in L^{i}(E_{n})$, $i = 1, 2$,
then $g_{1} - f_{1} = f_{2} - g_{2} \in L^{1} \cap L^{2}$. Since the two
definitions of the Fourier transform coincide on $L^{1} \cap L^{2}$ we have
$\hat{g}_{1} - \hat{f}_{1} = \hat{f}_{2} - \hat{g}_{2}$. Thus,
$\hat{f}_{1} + \hat{f}_{2} = \hat{g}_{1} + \hat{g}_{2}$. This shows that the
Fourier transform is well-defined on $L^{1}(E_{n}) + L^{2}(E_{n})$. Since the
latter is easily seen to contain all the spaces $L^{p}(E_{n})$,
$1 \leq p \leq 2$, the Fourier transform is defined for all
$f \in L^{p}(E_{n})$. It is easily checked that the Fourier inversion problem
can be solved in terms of the Abel or Gauss means in this case as well.
Similarly, Theorem~\ref{thm:c1-1-4} has the following extension:

\begin{theorem}\label{thm:c1-2-6}
If $f \in L^{1}(E_{n})$ and $g \in L^{p}(E_{n})$, $1 \leq p \leq 2$, then
$h = f * g$ belongs to $L^{p}(E_{n})$ (see Theorem~\ref{thm:c1-1-3}) and
\[
\hat{h}(x)=\hat{f}(x)\hat{g}(x)
\]
for almost every $x$.
\end{theorem}

In the next section we extend the definition of the Fourier transform much
further. We shall see that this extension is consistent with the one we have
just given when it is applied to the spaces $L^{p}(E_{n})$, $1 \leq p \leq 2$.

\section{The Class of Tempered Distributions}

The basic idea in the theory of distributions is to consider them as linear
functionals on some space of ``regular'' functions---the so-called ``testing
functions\index{testing functions@testing functions}.'' The space of testing functions is assumed to be well-behaved with
respect to the operations (differentiation, Fourier transformation,
convolution, translation, etc.) we have been studying, and this is then
reflected in the properties of distributions. We are naturally led to the
definition of such a space of testing functions by the following
considerations. Suppose we want these operations to be defined on a function
space, $\mathscr{S}$, and to preserve it. Then, it would certainly have to
consist of functions that are indefinitely differentiable; this, in view of
part (i) of \eqref{eq:c1-1-9}, indicates that each function of
$\mathscr{S}$, after being multiplied by a polynomial, must still be in
$\mathscr{S}$. We therefore make the following definition: The space
$\mathscr{S}$ of testing functions is defined to be the class of all those
$C^{\infty}$ functions $\varphi$ on $E_{n}$ (i.e., all the partial derivatives
of $\varphi$ exist and are continuous) such that
\begin{equation}\label{eq:c1-3-1}
\sup_{x\in E_{n}}\left|x^{\alpha}(D^{\beta}\varphi)(x)\right|<\infty
\end{equation}
for all $n$-tuples $\alpha = (\alpha_{1}, \ldots, \alpha_{n})$ and
$\beta = (\beta_{1}, \ldots, \beta_{n})$ of nonnegative integers (we are using
the notation that was used in \eqref{eq:c1-1-9}). For example
$\varphi(x) = e^{-\delta|x|^{2}}$, $\delta > 0$, belongs to $\mathscr{S}$; on
the other hand, $\varphi(x) = e^{-\delta|x|}$ fails to be differentiable at
the origin and, therefore, does not belong to $\mathscr{S}$. $\mathscr{S}$
contains the space $\mathscr{D}$ of all $C^{\infty}$ functions with compact
support.

It is not immediately clear that $\mathscr{D}$ is nonempty. To find a function
in $\mathscr{D}$, consider the function $f$ whose values are $f(t) = e^{-1/t}$
if $t > 0$ and $f(t) = 0$ when $t \leq 0$. Then $f \in C^{\infty}$, is bounded
and so are all its derivatives. Let $\varphi(t) = f(1 + t)f(1 - t)$; then
$\varphi(t) = e^{-2/(1 - t^{2})}$ if $|t| < 1$, is zero otherwise and clearly
belongs to $\mathscr{D} = \mathscr{D}(E_{1})$. We can easily obtain
$n$-dimensional variants from $\varphi$:

\textup{(a)} For $x \in E_{n}$ define
$\psi(x) = \varphi(x_{1})\varphi(x_{2}) \cdots \varphi(x_{n})$; then
$\psi \in \mathscr{D} = \mathscr{D}(E_{n})$;

\textup{(b)} For $x \in E_{n}$ define $\psi(x) = e^{-2/(1 - |x|^{2})}$, if
$|x| < 1$, and $0$ otherwise, then $\psi \in \mathscr{D} = \mathscr{D}(E_{n})$.

\textup{(c)} If $\eta \in C^{\infty}$ and $\psi$ is the function in (b) then
$\psi(\varepsilon x)\eta(x)$ defines a function in $\mathscr{D}(E_{n})$;
moreover $e^{2}\psi(\varepsilon x)\eta(x) \to \eta(x)$ as $\varepsilon \to 0$
(see \eqref{eq:c1-3-9} below).

We observe that the order of multiplication by powers of $x_{1}, \ldots, x_{n}$
and differentiation, in \eqref{eq:c1-3-1}, could have been reversed. That is,
$\varphi \in \mathscr{S}$ if and only if $\varphi \in C^{\infty}$ and
$\sup_{x \in E_{n}} |D^{\beta}(x^{\alpha}\varphi(x))| < \infty$ for all
$n$-tuples $\alpha$, $\beta$ of nonnegative integers. This shows that if $P$
is a polynomial in $n$ variables and $\varphi \in \mathscr{S}$ then
$P(x)\varphi(x)$ and $P(D)\varphi(x)$ are again in $\mathscr{S}$.

The spaces $C_{0}$ and $L^{p}(E_{n})$, $1 \leq p < \infty$, contain
$\mathscr{S}$. The latter is a dense subspace of each of the former; in fact,
the same is true of $\mathscr{D}$. The $L^{p}$ norm of $\varphi \in \mathscr{S}$
is bounded by a finite linear combination of $L^{\infty}$ norms of terms of
the form $x^{\alpha}\varphi(x)$. We see this in the following way: Let
$A = \|\varphi\|_{\infty}$ and
$B = \sup_{x \in E_{n}} |x|^{2n}|\varphi(x)|$ ($B < \infty$ by
\eqref{eq:c1-3-1}) then, for $1 \leq p < \infty$,
\begin{align*}
\left(\int_{E_{n}}|\varphi(x)|^{p}\,dx\right)^{1/p}
&\leq\left(\int_{|x|\leq1}|\varphi(x)|^{p}\,dx\right)^{1/p}
+\left(\int_{|x|>1}|\varphi(x)|^{p}\,dx\right)^{1/p}\\
&\leq A\left(\int_{|x|\leq1}1\,dx\right)^{1/p}
+B\left(\int_{|x|>1}|x|^{-2np}\,dx\right)^{1/p}\\
&=A\left(\frac{\omega_{n-1}}{n}\right)^{1/p}
+B\left(\frac{\omega_{n-1}}{(2p-1)n}\right)^{1/p}.
\end{align*}

This shows that the coefficients and exponents $\alpha$ involved in the linear
combination of $L^{\infty}$ norms majorizing $\|\varphi\|_{p}$ are independent
of $\varphi$.

From part (ii) of \eqref{eq:c1-1-9} one easily deduces that
$\hat{\varphi} \in C^{\infty}$ whenever $\varphi \in \mathscr{S}$. Combining
this with part (i), it then follows that $\hat{\varphi} \in \mathscr{S}$. That
is,

\begin{theorem}\label{thm:c1-3-2}
If $\varphi \in \mathscr{S}$ then $\hat{\varphi} \in \mathscr{S}$.
\end{theorem}

In fact, we see from the discussion on Fourier inversion in Section 1 that the
Fourier transform must be a one to one mapping of $\mathscr{S}$ onto
$\mathscr{S}$.

If $\varphi$ and $\psi$ belong to $\mathscr{S}$, then
Theorem~\ref{thm:c1-3-2} implies that $\hat{\varphi}$ and $\hat{\psi}$ belong
to $\mathscr{S}$ and, therefore, so does $\hat{\varphi}\hat{\psi}$. Since
$(\varphi * \psi)^{\wedge} = \hat{\varphi}\hat{\psi}$, an application of the
inverse Fourier transform shows that

\begin{theorem}\label{thm:c1-3-3}
If $\varphi$ and $\psi$ are in $\mathscr{S}$ so is $\varphi * \psi$.
\end{theorem}

We shall now introduce a metric on $\mathscr{S}$ that will make this space a
topological vector space. Toward this end we introduce a countable family of
norms, $\{\rho_{\alpha\beta}\}$, indexed by all the ordered pairs
$(\alpha, \beta)$ of $n$-tuples of nonnegative integers. For such a pair we
define, in view of \eqref{eq:c1-3-1},
\[
\rho_{\alpha\beta}(\varphi)=\sup_{x\in E_{n}}|x^{\alpha}D^{\beta}\varphi(x)|
\]
for $\varphi \in \mathscr{S}$. If we define
$d_{\alpha\beta}'(\varphi, \psi) = \rho_{\alpha\beta}(\varphi - \psi)$ then
$d_{\alpha\beta}'$ is a metric on $\mathscr{S}$. Let $d_{1}', d_{2}', \ldots$
be an ordering of these metrics and $d_{n} = d_{n}'/(1 + d_{n}')$,
$n = 1, 2, \ldots$. Then $d_{n}$ is a metric that is equivalent to $d_{n}'$
(that is, the two metrics define the same topology on $\mathscr{S}$).
Furthermore, $d_{n} \leq 1$. Therefore,
$d = \sum_{n=1}^{\infty} 2^{-n}d_{n}$ is a metric on $\mathscr{S}$. This is
the metric that will define the topology of $\mathscr{S}$. It is clear that
$\varphi_{k} \to \varphi$, with respect to $d$, if and only if $\varphi_{k} \to \varphi$,
with respect to each $d_{n}$ (as $k \to \infty$). From this it follows that
the vector space operations $(\varphi, \psi) \to \varphi + \psi$ and
$(a, \varphi) \to a\varphi$ ($a$ being a complex number) are continuous;
therefore, $(\mathscr{S}, d)$ is a topological vector space.

Some easily established properties of $\mathscr{S}$ and its topology, are the
following:

\refstepcounter{equation}\label{eq:c1-3-4}\textbf{(\theequation)}\quad
The mapping $\varphi(x) \to x^{\alpha}D^{\beta}\varphi(x)$ is continuous.

\refstepcounter{equation}\label{eq:c1-3-5}\textbf{(\theequation)}\quad
If $\varphi \in \mathscr{S}$ then $\lim_{h \to 0} \tau_{h}\varphi = \varphi$.

\refstepcounter{equation}\label{eq:c1-3-6}\textbf{(\theequation)}\quad
Suppose $\varphi \in \mathscr{S}$ and $h = (0, \ldots, h_{i}, \ldots, 0)$ lies
on the $i$-th coordinate axis of $E_{n}$ then the difference quotient
$[\varphi - \tau_{h}\varphi]/h_{i}$ tends to $\partial\varphi/\partial x_{i}$
as $|h| \to 0$.

\refstepcounter{equation}\label{eq:c1-3-7}\textbf{(\theequation)}\quad
$\mathscr{S}$ is a complete metric space.

\refstepcounter{equation}\label{eq:c1-3-8}\textbf{(\theequation)}\quad
The Fourier transform is a homeomorphism of $\mathscr{S}$ onto itself.

\refstepcounter{equation}\label{eq:c1-3-9}\textbf{(\theequation)}\quad
$\mathscr{D}$ is a dense subset of $\mathscr{S}$.

\refstepcounter{equation}\label{eq:c1-3-10}\textbf{(\theequation)}\quad
$\mathscr{S}$ is separable.

The collection $\mathscr{S}'$ of all continuous linear functionals on
$\mathscr{S}$ is called the space of tempered distributions\index{tempered distributions@tempered distributions}. We describe some
examples of tempered distributions:

\textup{(1)} Let $f \in L^{p}(E_{n})$, $1 \leq p \leq \infty$, and define
$L = L_{f}$ by letting
\[
L(\varphi)=L_{f}(\varphi)=\int_{E_{n}}f(x)\varphi(x)\,dx
\]
for $\varphi \in \mathscr{S}$. It is clear that $L_{f}$ is a linear functional
on $\mathscr{S}$. To show that it is continuous, therefore, it suffices to
show that it is continuous at the origin. Suppose then, $\varphi_{k} \to 0$
(in $\mathscr{S}$) as $k \to \infty$. We have seen that for any $q \geq 1$,
$\|\varphi_{k}\|_{q}$ is dominated by a finite linear combination of
$L^{\infty}$ norms of terms of the form $x^{\alpha}\varphi_{k}(x)$ (the
coefficients of this linear combination and the exponents $\alpha$ depend only
on $n$ and $q$ and not on $\varphi_{k}$). That is, $\|\varphi_{k}\|_{q}$ is
dominated by a finite linear combination of norms $\rho_{\alpha 0}$ evaluated
at $\varphi_{k}$. Thus $\|\varphi_{k}\|_{q} \to 0$ as $k \to \infty$.
Choosing $q$ so that $1/p + 1/q = 1$, H\"older's inequality shows
$|L(\varphi_{k})| \leq \|f\|_{p}\|\varphi_{k}\|_{q} \to 0$ as $k \to \infty$.
Thus, $L \in \mathscr{S}'$.

\textup{(2)} If $\mu$ is a finite Borel measure, the linear functional
$L = L_{\mu}$ defined by
\[
L(\varphi)=L_{\mu}(\varphi)=\int_{E_{n}}\varphi(x)\,d\mu(x)
\]
for $\varphi \in \mathscr{S}$ is a tempered distribution (choosing
$q = \infty$, the proof is the same as the one just given in the first
example).

\textup{(3)} A measurable function $f$ such that $f(x)/(1 + |x|^{2})^{k}$ is
in $L^{p}(E_{n})$, $1 \leq p \leq \infty$ (for some positive integer $k$), is
called a tempered\index{tempered distributions@tempered distributions!function or measure@function or measure} $L^{p}$ function (when $p = \infty$ such a function is often
also called a slowly increasing function\index{slowly increasing function@slowly increasing function}). For each such function,
$L(\varphi) = \int_{E_{n}} f(x)\varphi(x) \, dx$, $\varphi \in \mathscr{S}$,
defines a member of $\mathscr{S}'$. We see this by writing
\[
L(\varphi)=\int_{E_{n}}\{(1+|x|^{2})^{k}\varphi(x)\}\{f(x)/(1+|x|^{2})^{k}\}\,dx
\]
and observing that the mapping $\varphi(x) \to (1 + |x|^{2})^{k}\varphi(x)$ is
continuous in $\mathscr{S}$; the result then follows from example (1).

\textup{(4)} A tempered measure is a Borel measure $\mu$ such that
\[
\int_{E_{n}}(1+|x|^{2})^{-k}\,d|\mu|(x)<\infty
\]
for some integer $k$. Just as in example 3, we can show that
$L(\varphi) = \int_{E_{n}} \varphi(x) \, d\mu(x)$ defines a tempered
distribution.

\textup{(5)} Fix $x_{0} \in E_{n}$ and an $n$-tuple $\beta$ of nonnegative
integers. An immediate consequence of the continuity of the norm
$\rho_{\alpha\beta}$ in $\mathscr{S}$ is the fact that
$L(\varphi) = D^{\beta}\varphi(x_{0})$, for $\varphi \in \mathscr{S}$, defines
a tempered distribution. A special case is the Dirac\index{Dirac delta function@Dirac delta function} $\delta$-function:
$L(\varphi) = \varphi(0)$. This distribution, however, is also obtainable from
the tempered measure having mass $1$ concentrated at the origin (example (4)).
When $D^{\beta} = \partial/\partial x_{i}$ (that is,
$L(\varphi) = (\partial\varphi/\partial x_{i})(0)$) we obtain an example of a
tempered distribution that is not a special case of one of the previous four
types of distributions.

The tempered distributions of example (1) (or, more generally, example (3))
are called functions. Similarly, examples (2) and (4) define the distributions
that are called measures. We shall write, in these cases, $f$ and $\mu$
instead of $L_{f}$ and $L_{\mu}$. These functions and measures may be
considered as embedded in $\mathscr{S}'$. If we put on $\mathscr{S}'$ the
weakest topology such that the linear functionals $L \to L(\varphi)$
($\varphi \in \mathscr{S}$) are continuous, it is easy to see that the spaces
$L^{p}(E_{n})$, $1 \leq p \leq \infty$, are continuously embedded in
$\mathscr{S}'$. The same is true for the space of all finite Borel measures on
$E_{n}$ (which is a Banach space with the norm $\|\mu\| = \int_{E_{n}} d|\mu|$).

There exists a simple and important characterization of tempered
distributions:

\begin{theorem}\label{thm:c1-3-11}
A linear functional $L$ on $\mathscr{S}$ is a tempered distribution if and
only if there exists a constant $C > 0$ and integers $m$ and $l$ such that
\[
\left|L(\varphi)\right|\leq C\sum_{|\alpha|\leq l,|\beta|\leq m}\rho_{\alpha\beta}(\varphi)
\]
for all $\varphi \in \mathscr{S}$.
\end{theorem}

\begin{proof}
It is clear that the existence of $C$, $l$ and $m$ imply the continuity of $L$.

Suppose $L$ is continuous. It follows from the definition of the metric that a
basis for the neighborhoods of the origin in $\mathscr{S}$ is the collection
of sets
$N_{\varepsilon,l,m}=\{\varphi; \sum_{|\alpha|\leq l,|\beta|\leq m} \rho_{\alpha\beta}(\varphi) < \varepsilon\}$,
where $\varepsilon > 0$ and $l$ and $m$ are integers (for in the topology
induced by this system of neighborhoods---and their translates---$\varphi_{k} \to \varphi$
as $k \to \infty$ if and only if $\rho_{\alpha\beta}(\varphi_{k} - \varphi) \to 0$
for all $(\alpha, \beta)$). Thus, there exists such a set $N_{\varepsilon,l,m}$
such that $|L(\varphi)| \leq 1$ whenever $\varphi \in N_{\varepsilon,l,m}$. Let
$\|\varphi\| = \sum_{|\alpha|\leq l,|\beta|\leq m} \rho_{\alpha\beta}(\varphi)$
for all $\varphi \in \mathscr{S}$. If $0 < \bar{\varepsilon} < \varepsilon$
then $\psi = (\bar{\varepsilon}/\|\varphi\|)\varphi \in N_{\varepsilon,l,m}$
(if $\varphi \neq 0$). From the linearity of $L$ we obtain, therefore,
\[
(\bar{\varepsilon}/\|\varphi\|)|L(\varphi)|=|L(\psi)|\leq1.
\]
But this is the desired inequality with $C = 1/\bar{\varepsilon}$.
\end{proof}

We shall now show that some important operations in analysis
(differentiation, convolution, Fourier transformation, etc.) can be defined on
the class $\mathscr{S}'$. We begin by defining the convolution of a
distribution with a\index{convolution@convolution!with a distribution@with a distribution} testing function.

Toward this end we make the following definition: If $g$ is any function on
$E_{n}$ we define its reflection\index{reflection@reflection}, $\tilde{g}$, by letting
$\tilde{g}(x) = g(-x)$. A direct application of Fubini's theorem shows that if
$u$, $\varphi$, and $\psi$ are all in $\mathscr{S}$ then
\[
\int_{E_{n}}(u*\varphi)(x)\psi(x)\,dx=\int_{E_{n}}u(x)(\tilde{\varphi}*\psi)(x)\,dx.
\]
The mappings $\psi \to \int_{E_{n}}(u * \varphi)(x) \psi(x) \, dx$ and
$\theta \to \int_{E_{n}} u(x) \theta(x) \, dx$ are linear functionals on
$\mathscr{S}$. If we denote these functionals by $u * \varphi$ and $u$, the
last equality can be written in the form:
\begin{equation}\label{eq:c1-3-12}
(u*\varphi)(\psi)=u(\tilde{\varphi}*\psi).
\end{equation}
If $u \in \mathscr{S}'$ and $\varphi$, $\psi$ belong to $\mathscr{S}$, the
right side of \eqref{eq:c1-3-12} is well-defined since
$\tilde{\varphi} * \psi \in \mathscr{S}$. Furthermore, the mapping
$\psi \to u(\tilde{\varphi} * \psi)$, being the composition of two continuous
functions, is continuous. Thus, we can define the convolution of the
distribution $u$ with the testing function $\varphi$, $u * \varphi$, by means
of equality \eqref{eq:c1-3-12}.

It is easy to show that this convolution is associative in the sense that
$(u * \varphi) * \psi = u * (\varphi * \psi)$ whenever $u \in \mathscr{S}'$
and $\varphi, \psi \in \mathscr{S}$. The following result is a
characterization of the convolution we have just described.

\begin{theorem}\label{thm:c1-3-13}
If $u \in \mathscr{S}'$ and $\varphi \in \mathscr{S}$ then the convolution
$u * \varphi$ is the function $f$, whose value at $x \in E_{n}$ is
$f(x) = u(\tau_{x}\tilde{\varphi})$, where $\tau_{x}$ denotes the translation
by $x$ operator. Moreover, $f$ belongs to the class $C^{\infty}$ and it, as
well as all its derivatives, are slowly increasing.
\end{theorem}

\begin{proof}
We first show that $f$ is a $C^{\infty}$ slowly increasing function. Let
$h = (0, \ldots, h_{j}, \ldots, 0)$; then, by \eqref{eq:c1-3-6},
$[\tau_{x+h}\tilde{\varphi} - \tau_{x}\tilde{\varphi}]/h_{j} \to -\tau_{x}(\partial\tilde{\varphi}/\partial x_{j})$
as $|h| \to 0$, in the topology of $\mathscr{S}$. Thus, since $u$ is
continuous $[f(x + h) - f(x)]/h_{j} = u([\tau_{x+h}\tilde{\varphi} - \tau_{x}\tilde{\varphi}]/h_{j}) \to u[-\tau_{x}(\partial\tilde{\varphi}/\partial x_{j})]$
as $h_{j} \to 0$. This, together with \eqref{eq:c1-3-5}, shows that $f$ has
continuous first partial derivatives. Since
$\partial\tilde{\varphi}/\partial x_{j} \in \mathscr{S}$ we can iterate this
argument and show that $D^{\beta}f$ exists and is continuous for all
$n$-tuples of nonnegative integers $\beta$. We observe that we obtain
$(D^{\beta}f)(x) = (-1)^{|\beta|} u(\tau_{x}D^{\beta}\tilde{\varphi})$, where
$|\beta| = \beta_{1} + \cdots + \beta_{n}$. Consequently, since
$D^{\beta}\tilde{\varphi} \in \mathscr{S}$ if $f$ were slowly increasing, the
same would hold for all the derivatives of $f$. The fact that $f$ is slowly
increasing, however, is an easy consequence of Theorem~\ref{thm:c1-3-11}:
There exist $C > 0$ and integers $m$ and $l$ such that
\[
\left|f(x)\right|=\left|u(\tau_{x}\tilde{\varphi})\right|
\leq C\sum_{|\alpha|\leq l,|\beta|\leq m}\rho_{\alpha\beta}(\tau_{x}\tilde{\varphi}).
\]
But
$\rho_{\alpha\beta}(\tau_{x}\tilde{\varphi}) = \sup_{w \in E_{n}} |w^{\alpha}(D^{\beta}\tilde{\varphi})(w - x)| = \sup_{w \in E_{n}} |(w + x)^{\alpha}(D^{\beta}\tilde{\varphi})(w)|$
and the latter is clearly bounded by a polynomial in $x$.

In order to show that $u * \varphi$ is the function $f$ we must show that
$(u * \varphi)(\psi) = \int_{E_{n}} \psi(t)f(t) \, dt$. But,
\begin{align*}
(u*\varphi)(\psi)=u(\tilde{\varphi}*\psi)
&=u\left(\int_{E_{n}}\tilde{\varphi}(x-t)\psi(t)\,dt\right)\\
&=u\left(\int_{E_{n}}(\tau_{t}\tilde{\varphi})(x)\psi(t)\,dt\right).
\end{align*}
On the other hand, the Riemann sums of the last integral are easily shown to
converge in the topology of $\mathscr{S}$. Thus,
\[
u\left(\int_{E_{n}}(\tau_{t}\tilde{\varphi})(x)\psi(t)\,dt\right)
=\int_{E_{n}}u(\tau_{t}\tilde{\varphi})\psi(t)\,dt
=\int_{E_{n}}f(t)\psi(t)\,dt,
\]
since $u$ is continuous and linear, which is the desired equality.
\end{proof}

We now examine differentiation in the class $\mathscr{S}'$. We first observe
that integration by parts implies
\[
\int_{E_{n}}(D^{\beta}u)(x)\varphi(x)\,dx=(-1)^{|\beta|}\int_{E_{n}}u(x)(D^{\beta}\varphi)(x)\,dx
\]
for $u, \varphi \in \mathscr{S}$. The mappings
$\varphi \to \int_{E_{n}} (D^{\beta}u)(x)\varphi(x) \, dx$ and
\[
\psi\to\int_{E_{n}}u(x)\psi(x)\,dx
\]
are continuous linear functionals on $\mathscr{S}$. Denoting them by
$(D^{\beta}u)$ and $u$, the last equality can be written in the form
\begin{equation}\label{eq:c1-3-14}
(D^{\beta}u)(\varphi)=(-1)^{|\beta|}u(D^{\beta}\varphi).
\end{equation}
But the right side is well-defined whenever $u \in \mathscr{S}'$ and
$\varphi \in \mathscr{S}$. Furthermore, $\varphi \to u(D^{\beta}\varphi)$ is
continuous on $\mathscr{S}$ (being the composite of two continuous
functions). Therefore, we can define the partial derivative\index{derivatives@derivatives!of distributions@of distributions} $D^{\beta}u$ of
the distribution $u$ by means of equality \eqref{eq:c1-3-14}. It is then clear
that $D^{\beta}u$ is a member of $\mathscr{S}'$.

In a similar way, by first considering $u$, $\varphi \in \mathscr{S}$, we can
motivate the following two definitions: The translation operator $\tau_{h}$ is
defined on $\mathscr{S}'$ by letting $\tau_{h}u$, when $u$ belongs to
$\mathscr{S}'$, be the continuous linear functional on $\mathscr{S}$ whose
value at $\varphi$ is given by $(\tau_{h}u)(\varphi) = u(\tau_{-h}\varphi)$.
Putting $\tilde{u}(\varphi) = u(\tilde{\varphi})$ for all
$\varphi \in \mathscr{S}$ we obtain the reflection, $\tilde{u}$, of the
distribution $u$.

By the multiplication formula (Theorem~\ref{thm:c1-1-15}) we have
$u(\hat{\varphi}) = \int_{E_{n}} u(x)\hat{\varphi}(x) \, dx = \int_{E_{n}} \hat{u}(x)\varphi(x) \, dx = \hat{u}(\varphi)$
for all $u, \varphi \in \mathscr{S}$. Therefore, we define the Fourier
transform\index{Fourier transform@Fourier transform!of distributions@of distributions} $\hat{u}$ of the distribution $u$ to be the continuous linear
functional whose value as $\varphi \in \mathscr{S}$ is
\begin{equation}\label{eq:c1-3-15}
\hat{u}(\varphi)=u(\hat{\varphi}).
\end{equation}
We observe that if $f$ belongs to $L^{p}(E_{n})$, $1 \leq p \leq 2$, its
Fourier transform as a distribution coincides with the function $\hat{f}$
defined at the end of the second section.%
\footnote{For the case $p > 2$, see \eqref{eq:c1-4-13} below.}
It is easy to check that the Fourier transform we have just defined is an
isomorphism of the topological vector space $\mathscr{S}'$ onto itself.

Having set down these facts of distribution theory we shall now apply them to
the study of the basic class of linear operators that occur in Fourier
analysis: the class of operators that commute\index{commuting with translations (operators)@commuting with translations (operators)} with translations. While the
problem of characterizing such operators in the general context of $L^{p}$
spaces remains open, it can be solved in certain significant special cases, as
we shall see below.

Suppose $B$ is an operator mapping a linear space $V$ of functions on $E_{n}$
into another such space $W$. We say that $B$ commutes with translations if
$\tau_{h}B = B\tau_{h}$ for all $h \in E_{n}$. An example of such an operator
is obtained as follows: Let us fix an $f \in L^{p}(E_{n})$ and define
$Bg = f * g$ for $g \in L^{1}(E_{n})$. By Theorem~\ref{thm:c1-1-3} we then see
that $B$ is a bounded operator mapping $L^{1}(E_{n})$ into $L^{p}(E_{n})$ (we
have $\|B\| \leq \|f\|_{p}$) and a change of variables shows that
$\tau_{h}(Bg) = B(\tau_{h}g)$ for all $g \in L^{1}(E_{n})$. The following
result shows that ``all'' bounded operators commuting with translations are of
similar ``convolution type'':

\begin{theorem}\label{thm:c1-3-16}
Suppose $B\colon L^{p}(E_{n})\to L^{q}(E_{n})$, $1\leq p,q\leq\infty$, is
linear, bounded and commutes with translations; then there exists a unique
tempered distribution $u$ such that $B\varphi = u * \varphi$ for all
$\varphi \in \mathscr{S}$.%
\footnote{Consider, for simplicity, the case $n = 1$. By taking appropriate
Riemann sums we see that if $u$ is a well-behaved function, then
$(B_{\varphi})(x)$ is approximated by a linear combination of the form
$\sum_{k=1}^{n} u(t_{k})(t_{k}-t_{k-1})\varphi(x-t_{k})=\sum_{k=1}^{n} c_{k}\varphi(x-t_{k})$.
That is, $B$ is approximated by a finite linear combination of translation
operators. In a sense that is made precise by Theorem~\ref{thm:c1-3-16} all
bounded linear operators commuting with translations can be so approximated.}
\end{theorem}

We shall show that the proof of this theorem is an easy consequence of the
following lemma:

\begin{lemma}\label{lem:c1-3-17}
If $f \in L^{p}(E_{n})$ has derivatives in the $L^{p}$ norm of all orders
$\leq n + 1$, then $f$ equals almost everywhere a continuous function $g$
satisfying
\[
\left|g(0)\right|\leq C\sum_{|\alpha|\leq n+1}\left\|D^{\alpha}f\right\|_{p},
\]
where $C$ depends only on the dimension $n$ and the exponent $p$.
\end{lemma}

\begin{proof}
Let $x = (x_{1}, \ldots, x_{n}) \in E_{n}$. Then there exists
$C' = C_{n}'$ such that
\[
(1+|x|^{2})^{(n+1)/2}\leq(1+|x_{1}|+\cdots+|x_{n}|)^{n+1}\leq C^{\prime}\sum_{|\alpha|\leq n+1}|x^{\alpha}|.
\]
Let us first suppose $p = 1$. Then by part (ii) of \eqref{eq:c1-1-9}, and
Theorem~\ref{thm:c1-1-1} (a)
\begin{align*}
|\hat{f}(x)|
&\leq C^{\prime}(1+|x|^{2})^{-[(n+1)/2]}\sum_{|\alpha|\leq n+1}|x^{\alpha}||\hat{f}(x)|\\
&=C^{\prime}(1+|x|^{2})^{-[(n+1)/2]}\sum_{|\alpha|\leq n+1}|((2\pi)^{-|\alpha|}D^{\alpha}f)^{\wedge}(x)|\\
&\leq C^{\prime\prime}(1+|x|^{2})^{-[(n+1)/2]}\sum_{|\alpha|\leq n+1}\|D^{\alpha}f\|_{1}.
\end{align*}
Since $(1 + |x|^{2})^{-[(n+1)/2]}$ defines an integrable function on $E_{n}$,
it follows that $\hat{f} \in L^{1}(E_{n})$ and, letting
$C = C'' \int_{E_{n}} (1 + |x|^{2})^{-[(n+1)/2]} \, dx$, we have
\[
\left\|\hat{f}\right\|_{1}\leq C\sum_{|\alpha|\leq n+1}\left\|D^{\alpha}f\right\|_{1}.
\]
Thus, by Corollary~\ref{cor:c1-1-21} $f$ equals almost everywhere a continuous
function $g$ and, by Theorem~\ref{thm:c1-1-1} (a),
\[
\left|g(0)\right|\leq\left\|f\right\|_{\infty}\leq\left\|\hat{f}\right\|_{1}
\leq C\sum_{|\alpha|\leq n+1}\left\|D^{\alpha}f\right\|_{1}.
\]
Suppose now that $p > 1$. Choose $\varphi \in C^{\infty}$ such that
$\varphi(x) = 1$ if $|x| \leq 1$ and $\varphi(x) = 0$ if $|x| > 2$. Then
$\varphi f$ satisfies the conditions of the lemma for $p = 1$. Thus,
$\varphi f$ equals almost everywhere a continuous function $h$ such that
\[
\left|h(0)\right|\leq C\sum_{|\alpha|\leq n+1}\left\|D^{\alpha}(\varphi f)\right\|_{1}.
\]
Since $D^{\alpha}(\varphi f) = \sum_{\mu + \nu = \alpha} C^{\mu\nu}(D^{\mu}f)(D^{\nu}\varphi)$,
we have
\begin{align*}
\|D^{\alpha}(\varphi f)\|_{1}
&\leq\int_{|x|\leq2}\sum_{\mu+\nu=\alpha}C^{\mu\nu}|D^{\mu}f||D^{\nu}\varphi|\,dx\\
&\leq\sum_{\mu+\nu=\alpha}C\Biggl\{\sup_{|x|\leq2}|(D^{\nu}\varphi)(x)|\Biggr\}\int_{|x|\leq2}|D^{\mu}f(x)|\,dx\\
&\leq A\sum_{|\nu|\leq|\alpha|}\int_{|x|\leq2}|D^{\nu}f(x)|\,dx
\leq AB\sum_{|\nu|\leq|\alpha|}\|D^{\nu}f\|_{p},
\end{align*}
where $A \geq \|D^{\nu}\varphi\|_{\infty}$, $|\nu| \leq |\alpha|$, and $B$
depends only on $p$ and $n$. Thus, we can find a constant $K$ such that
\[
\left|h(0)\right|\leq K\sum_{|\alpha|\leq n+1}\left\|D^{\alpha}f\right\|_{p}.
\]
Since $\varphi(x) = 1$ if $|x| \leq 1$ we see that $f$ is equal almost
everywhere to a continuous function $g$ in the sphere of radius $1$ about $0$;
moreover,
\[
\left|g(0)\right|=\left|h(0)\right|\leq K\sum_{|\alpha|\leq n+1}\left\|D^{\alpha}f\right\|_{p}.
\]
But, by choosing $\varphi$ appropriately, the argument clearly shows that $f$
equals almost everywhere a continuous function on any sphere about $0$. This
proves the lemma.
\end{proof}

We now turn to the proof of Theorem~\ref{thm:c1-3-16}. We first observe that
if $\varphi \in \mathscr{S}$ then $B\varphi$ has derivatives in the $L^{q}$
norm of all orders. For, if $h = (0, \ldots, h_{j}, \ldots, 0)$ lies on the
$j$-th coordinate axis we have
$(\tau_{h}(B\varphi) - B\varphi)/h_{j} = (B(\tau_{h}\varphi) - B\varphi)/h_{j} = B((\tau_{h}\varphi - \varphi)/h_{j})$.
Since $(\tau_{h}\varphi - \varphi)/h_{j}$ converges in $\mathscr{S}$ to
$-\varphi_{j} = -\partial\varphi/\partial x_{j}$, it converges to this
function in the $L^{p}$ norm. Since $B$ is a bounded operator from $L^{p}$ to
$L^{q}$ it follows, therefore, that
$(\tau_{h}(B\varphi) - B\varphi)/h_{j}$ converges to
$-\partial(B\varphi)/\partial x_{j} = -B\varphi_{j}$ in the $L^{q}$ norm. The
fact that $B\varphi$ has derivatives in the $L^{q}$ norm of all orders now
follows by iteration. Clearly, $B(D^{\alpha}\varphi) = D^{\alpha}(B\varphi)$
for all $n$-tuples $\alpha = (\alpha_{1}, \ldots, \alpha_{n})$ of nonnegative
integers. By Lemma~\ref{lem:c1-3-17}, therefore, $B\varphi$ equals almost
everywhere a continuous function $g_{\varphi}$ satisfying
\begin{align*}
|g_{\varphi}(0)|
&\leq C\sum_{|\alpha|\leq n+1}\|D^{\alpha}(B\varphi)\|_{q}\\
&=C\sum_{|\alpha|\leq n+1}\|B(D^{\alpha}\varphi)\|_{q}
\leq\|B\|C\sum_{|\alpha|\leq n+1}\|D^{\alpha}\varphi\|_{p}.
\end{align*}

This last inequality shows that the mapping $\varphi \to g_{\varphi}(0)$ is a
continuous linear functional, $u_{1}$, on $\mathscr{S}$ (this follows from
Theorem~\ref{thm:c1-3-11} and the observation, made at the beginning of this
section, that the $L^{p}$ norm of a $\psi \in \mathscr{S}$ is bounded by a
finite linear combination of $L^{\infty}$ norms of terms of the form
$x^{\alpha}\psi(x)$). We claim that $u = \tilde{u}_{1}$ is the linear
functional we are seeking. Indeed, if $\varphi \in \mathscr{S}$, using
Theorem~\ref{thm:c1-3-13} we obtain
$(u * \varphi)(x) = u(\tau_{x}\tilde{\varphi}) = u([\tau_{-x}\varphi]^{\sim}) = \tilde{u}(\tau_{-x}\varphi) = u_{1}(\tau_{-x}\varphi) = (B(\tau_{-x}\varphi))(0) = (\tau_{-x}B\varphi)(0) = (B\varphi)(x)$.%
\footnote{We are identifying $B(\tau_{-x}\varphi)$ with the continuous
function $g_{\tau_{-x}\varphi}$ that it equals a.e.}

We note that it follows from this construction that $u$ is unique. The
theorem is therefore proved.

Combining this result with Theorem~\ref{thm:c1-3-13} we obtain the fact that
$B\varphi$, for $\varphi \in \mathscr{S}$, is almost everywhere equal to a
$C^{\infty}$ function which, together with all its derivatives, is slowly
increasing.

Let $(L^{p}, L^{q})$ be the collection of those tempered distributions $u$ for
which there exists an $A > 0$ such that $\|u * \varphi\|_{q} \leq A \|\varphi\|_{p}$
for all $\varphi \in \mathscr{S}$. Theorem~\ref{thm:c1-3-16} when
$p < \infty$ gives a one-to-one correspondence between this set and the
bounded linear operators from $L^{p}$ into $L^{q}$ that commute with
translations (that the operator $\varphi \to u * \varphi$ commutes with
translations can be easily seen by using Theorem~\ref{thm:c1-3-13}). When
$p = 2 = q$ one can give a very simple characterization of this set of
distributions:

\begin{theorem}\label{thm:c1-3-18}
The distribution $u$ belongs to $(L^{2}, L^{2})$ if and only if there exists
$b \in L^{\infty}(E_{n})$ such that $\hat{u} = b$. In this case,
$\|b\|_{\infty}$ is the norm of the operator
$B\colon L^{2} \cap \mathscr{S} \to L^{2}$ defined by putting
$B\varphi = u * \varphi$; moreover, $(u * \varphi)^{\wedge} = \hat{u}\hat{\varphi}$.
\end{theorem}

\begin{proof}
If $v \in \mathscr{S}'$ and $\psi \in \mathscr{S}$ we define their product,
$v\psi$, to be the element of $\mathscr{S}'$ such that $(v\psi)(\varphi) = v(\psi\varphi)$
for all $\varphi \in \mathscr{S}$. With the product\index{product@product!of distribution and testing function@of distribution and testing function} of a distribution with a
testing function so defined we first observe that whenever $u \in \mathscr{S}'$
and $\varphi \in \mathscr{S}$ then
\[
\textup{(i)}\qquad (u*\varphi)^{\wedge}=\hat{u}\hat{\varphi}.
\]
To see this we must show that
$(u * \varphi)^{\wedge}(\psi) = (\hat{u}\hat{\varphi})(\psi)$ for all
$\psi \in \mathscr{S}$. It follows immediately from
Theorem~\ref{thm:c1-1-4} and the Fourier inversion formula that
$\hat{\varphi}\psi$ is the inverse Fourier transform of
$\tilde{\varphi} * \hat{\psi}$. Thus, by the above definitions of the Fourier
transform and of convolution,
$(u * \varphi)^{\wedge}(\psi) = (u * \varphi)(\hat{\psi}) = u(\tilde{\varphi} * \hat{\psi}) = \hat{u}(\hat{\varphi}\psi) = (\hat{u}\hat{\varphi})(\psi)$;
hence, equality (i) is established.

Let $\varphi_{0} = e^{-\pi|x|^{2}}$; then $\varphi_{0} \in \mathscr{S}$ and
$\hat{\varphi}_{0} = \varphi_{0}$ (see Theorem~\ref{thm:c1-1-13} and footnote
4). If $u \in (L^{2}, L^{2})$ then, by the Plancherel theorem,
$\Phi_{0} = \hat{u}\varphi_{0} = (u * \varphi_{0})^{\wedge} \in L^{2}(E_{n})$.
Let $b(x) = e^{\pi|x|^{2}}\Phi_{0}(x) = \Phi_{0}(x)/\varphi_{0}(x)$.

We claim that
\[
\textup{(ii)}\qquad (u*\varphi)^{\wedge}=b\hat{\varphi}
\]
for all $\varphi \in \mathscr{S}$. By (i) it suffices to show that
$(\hat{u}\hat{\varphi})(\psi) = (b\hat{\varphi})(\psi)$ for all
$\psi \in \mathscr{D}$ (since $\mathscr{D}$ is dense in $\mathscr{S}$). But, if
$\psi \in \mathscr{D}$, then
$(\psi/\hat{\varphi}_{0})(x) = \psi(x)e^{\pi|x|^{2}} \in \mathscr{D}$; thus,
\begin{align*}
(\hat{u}\hat{\varphi})(\psi)
&=\hat{u}(\hat{\varphi}\psi)=\hat{u}(\hat{\varphi}\hat{\varphi}_{0}\psi/\hat{\varphi}_{0})
=(\hat{u}\hat{\varphi}_{0})(\hat{\varphi}\psi/\hat{\varphi}_{0})\\
&=\int_{E_{n}}\Phi_{0}(x)\hat{\varphi}(x)e^{\pi|x|^{2}}\psi(x)\,dx\\
&=\int_{E_{n}}b(x)\hat{\varphi}(x)\psi(x)\,dx=(b\hat{\varphi})(\psi).
\end{align*}

It follows immediately that $\hat{u} = b$: We have just shown that
$\hat{u}(\hat{\varphi}\psi) = (b\hat{\varphi})(\psi) = b(\hat{\varphi}\psi)$
for all $\varphi \in \mathscr{S}$ and $\psi \in \mathscr{D}$. Selecting
$\varphi$ so that $\hat{\varphi}(x) = 1$ for $x$ in the support of $\psi$,
this shows that $\hat{u}(\psi) = b(\psi)$ for all $\psi \in \mathscr{D}$.
Thus, $\hat{u} = b$.

Since $u \in (L^{2}, L^{2})$, there exists $A > 0$ such that
$\|b\hat{\varphi}\|_{2} = \|(u * \varphi)^{\wedge}\|_{2} = \|u * \varphi\|_{2} \leq A\|\varphi\|_{2}$
for all $\varphi \in \mathscr{S}$. But this clearly implies that
$b \in L^{\infty}(E_{n})$ with $\|b\|_{\infty} \leq A$ (since $\mathscr{S}$ is
dense in $L^{2}(E_{n})$).

On the other hand, if $\hat{u} = b \in L^{\infty}(E_{n})$ the Plancherel
theorem and equality (i) immediately imply that $u \in (L^{2}, L^{2})$ and
$\|b\|_{\infty}$ is the operator norm of $B$.
\end{proof}

Another simple characterization of the class $(L^{p}, L^{q})$ can be given
when $p = 1 = q$:

\begin{theorem}\label{thm:c1-3-19}
The distribution $u$ belongs to $(L^{1}, L^{1})$ if and only if it is a finite
Borel measure. In this case, the total variation of $u$ equals the norm of the
operator $B\colon L^{1} \cap \mathscr{S} \to L^{1}$ defined by
$B\varphi = u * \varphi$ for $\varphi \in \mathscr{S}$.
\end{theorem}

\begin{proof}
If $u$ is a finite Borel measure it is clear that it belongs to
$(L^{1}, L^{1})$ (see Remark following Theorem~\ref{thm:c1-1-3}).

On the other hand, if $u \in (L^{1}, L^{1})$ consider the family of $L^{1}$
functions $u_{\varepsilon} = (u * W(\cdot, \varepsilon))$, $\varepsilon > 0$
(the Gauss--Weierstrass kernel $W(\cdot, \varepsilon)$ belongs to
$\mathscr{S}$ for each $\varepsilon > 0$). Then, by our assumption on $u$ and
Lemma~\ref{lem:c1-1-17}, part (a),
$\|u_{\varepsilon}\|_{1} \leq A\|W(\cdot, \varepsilon)\|_{1} = A$. That is,
the family $\{u_{\varepsilon}\}$ is uniformly bounded in the $L^{1}$ norm. Let
us consider $L^{1}(E_{n})$ as embedded in the Banach space $M = M(E_{n})$ of
finite Borel measures on $E_{n}$.%
\footnote{If $f \in L^{1}(E_{n})$ then $\mu(F) = \int_{F} f(x) \, dx$, $F$ a
measurable subset of $E_{n}$, defines a finite Borel measure on $E_{n}$. The
total variation of $\mu$ equals $\|f\|_{1}$.}
$M(E_{n})$ can be identified with the dual of $C_{0} = C_{0}(E_{n})$ by making
each $\mu \in M$ correspond to the linear functional assigning to
$\varphi \in C_{0}$ the value $\int_{E_{n}} \varphi(x) \, d\mu(x)$. Thus, the
unit sphere of $M$ is compact in the weak$*$ topology.%
\footnote{That is, the weakest topology for which the members, $\varphi$, of
$C_{0}$ can be considered as continuous linear functionals of $M$ (if
$\mu \in M$ is mapped onto $\int_{E_{n}} \varphi(x) \, d\mu(x)$ we obtain a
linear functional on $M$). See bibliographical notes for details.}
In particular, we can find $\mu \in M$ and a null sequence
$\{\varepsilon_{k}\}$ such that $u_{\varepsilon_{k}} \to \mu$ as
$k \to \infty$ in this topology. That is, for each $\varphi \in C_{0}$,
\[
\textup{(i)}\qquad
\lim_{k\to\infty}\int_{E_{n}}\varphi(x)u_{\varepsilon_{k}}(x)\,dx
=\int_{E_{n}}\varphi(x)\,d\mu(x).
\]
We now claim that $\mu$, considered as a distribution, equals $u$.

We must show, therefore, that $u(\psi) = \int_{E_{n}} \psi(x) \, d\mu(x)$ for
all $\psi \in \mathscr{S}$. Let
$\psi_{\varepsilon}(x) = \int_{E_{n}} \psi(x - t)W(t, \varepsilon) \, dt$.
Then, for all $n$-tuples $\alpha = (\alpha_{1}, \ldots, \alpha_{n})$ of
nonnegative integers, we have
$(D^{\alpha}\psi_{\varepsilon})(x) = \int_{E_{n}} (D^{\alpha}\psi)(x - t)W(t, \varepsilon) \, dt$.
It follows from Theorem~\ref{thm:c1-1-18} that $(D^{\alpha}\psi_{\varepsilon})(x)$
converges to $(D^{\alpha}\psi)(x)$ uniformly in $x$. Thus,
$\psi_{\varepsilon} \to \psi$ as $\varepsilon \to 0$ (in $\mathscr{S}$) and
this implies that $u(\psi_{\varepsilon}) \to u(\psi)$. But, since
$W(\cdot, \varepsilon) = \tilde{W}(\cdot, \varepsilon)$,
\[
u(\psi_{\varepsilon})=u(W(\cdot,\varepsilon)*\psi)
=(u*W(\cdot,\varepsilon))(\psi)=\int_{E_{n}}\psi(x)u_{\varepsilon}(x)\,dx.
\]
Thus, putting $\varepsilon = \varepsilon_{k}$, letting $k \to \infty$ and
applying (i) with $\varphi = \psi$ we obtain the desired equality
$u(\psi) = \int_{E_{n}} \psi(x) \, d\mu(x)$. The remainder of the theorem now
follows easily.
\end{proof}

We mentioned before that a simple characterization of the general
$(L^{p}, L^{q})$ spaces is not known (and probably does not exist). We do
have, however, the following general result of duality:

\begin{theorem}\label{thm:c1-3-20}
If $1 \leq p$, $q \leq \infty$, $1/p + 1/p' = 1$ and $1/q + 1/q' = 1$ then
$(L^{p}, L^{q}) = (L^{q'}, L^{p'})$.
\end{theorem}

\begin{proof}
By the F. Riesz representation theorem we can identify the spaces $L^{p'}$ and
$L^{q'}$ with the duals of $L^{p}$ and $L^{q}$ (this is true if $p, q < \infty$.
If one of $p$ or $q$ is $\infty$, a slight modification of the argument we
give will prove the theorem in this case as well). Let $u \in (L^{p}, L^{q})$
and $B\colon L^{p} \to L^{q}$ denote the unique bounded linear extension to
$L^{p}$ of the mapping $\varphi \to u * \varphi$
($\varphi \in \mathscr{S}$). If $B^{*}\colon L^{q'} \to L^{p'}$ denotes the
adjoint to $B$ we then have
\[
\int_{E_{n}}(B\varphi)\psi\,dx=\int_{E_{n}}\varphi(B^{*}\psi)\,dx
\]
for all $\varphi$, $\psi \in \mathscr{S}$. But, in terms of distributions,
this is equivalent to the equality $(u * \varphi)(\psi) = (B^{*}\psi)(\varphi)$.
That is,
$(B^{*}\psi)(\varphi) = u(\tilde{\varphi} * \psi) = u([\tilde{\psi} * \varphi]^{\sim}) = \tilde{u}(\tilde{\psi} * \varphi) = (\tilde{u} * \psi)(\varphi)$
for all $\varphi$, $\psi \in \mathscr{S}$. Thus, $B^{*}\psi = \tilde{u} * \psi$
for all $\psi \in \mathscr{S}$. Since
$\|B^{*}\psi\|_{p'} \leq \|B^{*}\|\|\psi\|_{q'} = \|B\|\|\psi\|_{q'}$ we have
$\tilde{u} \in (L^{q'}, L^{p'})$. Since $u \in (L^{q'}, L^{p'})$ if and only if
$\tilde{u} \in (L^{q'}, L^{p'})$ this shows that
$(L^{p}, L^{q}) \subset (L^{q'}, L^{p'})$. The opposite inclusion relation
follows by interchanging the roles of $(p, q)$ and $(q', p')$.
\end{proof}

\section{Further Results}

\refstepcounter{equation}\label{eq:c1-4-1}\textbf{\theequation.}~We remarked,
after the proof of Theorem~\ref{thm:c1-1-2}, that belonging to $C_{0}$ is not a
sufficient condition for being the Fourier transform of an integrable
function. Suppose, for simplicity, that $n = 1$; then one way of seeing this
is the following: We first observe that if $\hat{f}$ is such a Fourier
transform and, moreover, $\hat{f}$ is an odd function then, for
$1 < b < \infty$, $|\int_{1}^{b} [\hat{f}(x)/x] \, dx| \leq A$, where $A$ does
not depend on $b$. This is a consequence of the well-known corresponding
property for the integrals of $\sin x/x$; that is
$|\int_{\alpha}^{\beta} (\sin x/x) \, dx| \leq B < \infty$ independently of
$\alpha$, $\beta$ ($0 \leq |\alpha| < |\beta| < \infty$). Indeed, because of
the oddness of $\hat{f}$ we have
$\hat{f}(x) = -i \int_{-\infty}^{\infty} f(t) \sin (2\pi xt) \, dt$. Thus,
using Fubini's theorem one easily deduces that
$|\int_{1}^{b} [\hat{f}(x)/x] \, dx| \leq A$. To give an example, therefore, of
a function in $C_{0}$ that is not a Fourier transform of an integrable
function all we need to do is exhibit a continuous odd function $g$ vanishing
at $\infty$ and such that $\int_{1}^{b} (g(x)/x) \, dx$ is not bounded for
$b \to \infty$. But this is clearly the case if, for example,
$g(x) = 1/\log x$ for $x$ large.

\refstepcounter{equation}\label{eq:c1-4-2}\textbf{\theequation.}~For
$g \in L^{1}(E_{n})$ fixed let us consider the operator $B$ mapping
$f \in L^{p}(E_{n})$ to $f * g$. From Theorem~\ref{thm:c1-1-3} we see that $B$
is a bounded linear transformation from $L^{p}(E_{n})$ into $L^{p}(E_{n})$
whose norm, $\|B\|^{(p)}$, is majorized by $\|g\|_{1}$. When $p = 1$ it is easy
to show that $\|B\|^{(1)} = \|g\|_{1}$ (for, by Theorem~\ref{thm:c1-1-18} we
see that $f_{\varepsilon}(t) = P(t, \varepsilon)$ is a family of functions of
$L^{1}$ norm equal to $1$ such that $f_{\varepsilon} * g \to g$ in the $L^{1}$
norm as $\varepsilon \to 0$). We have seen (Theorem~\ref{thm:c1-3-18}) that
$\|B\|^{(2)} = \|\hat{g}\|_{\infty}$. It is natural to ask, therefore, how the
norm $\|B\|^{(p)}$ can be expressed in terms of $g$ for $1 \leq p \leq \infty$.
A satisfactory answer is not known in general. If $g \geq 0$, however, then
$\|g\|_{1} = |\hat{g}(0)| \leq \|\hat{g}\|_{\infty} \leq \|g\|_{1}$; thus,
$\|B\|^{(2)} = \|g\|_{1}$. From this, Theorem~\ref{thm:c1-3-20} and a result
(the M. Riesz convexity theorem) in the theory of interpolation of operators
(that we shall develop in Chapter V) it follows that $\|B\|^{(p)} = \|g\|_{1}$,
for $1 \leq p \leq \infty$, provided $g \geq 0$.

\refstepcounter{equation}\label{eq:c1-4-3}\textbf{\theequation.}~Theorem~\ref{thm:c1-1-3}
has the following extension: if $f \in L^{p}(E_{n})$ and $g \in L^{r}(E_{n})$,
$1 \leq p, r$ and $(1/p) + (1/r) \geq 1$ then $h = f * g \in L^{q}(E_{n})$,
where $(1/q) = (1/p) + (1/r) - 1$, and
\[
\left\|h\right\|_{q}\leq\left\|f\right\|_{p}\left\|g\right\|_{r}.
\]
For a direct proof of this result (often called Young's inequality) see
Zygmund~\cite{Zygmund1}, Chapter II, page 37. We shall show, in Chapter V, that
this inequality also follows immediately from the M. Riesz convexity theorem.

\refstepcounter{equation}\label{eq:c1-4-4}\textbf{\theequation.}~We can obtain
more insight into the meaning of the derivatives in the $L^{p}$ norms by
observing that (a) when the dimension $n$ is 1, then $f \in L^{p}(E_{1})$ has
a derivative in the $L^{p}$ norm if and only if $f$ is equal almost everywhere
to a locally absolutely continuous function whose derivative, $f'$, belongs to
$L^{p}(E_{1})$ (see Bochner and Chandrasekharan~\cite{BochnerChandrasekharan1},
Section 6); (b) in general, $f \in L^{p}(E_{n})$ has a $k$-th partial
derivative in the $L^{p}$ norm if and only if $f$, regarded as a tempered
distribution, has a $k$-th partial derivative, $(\partial/\partial x_{k})f$, in
the distribution sense (see \eqref{eq:c1-3-14}) that is an $L^{p}$ function.

\refstepcounter{equation}\label{eq:c1-4-5}\textbf{\theequation.}~If
$f \in L^{p}(E_{n})$, $1 < p \leq \infty$, is regarded as a tempered
distribution then all its first partial derivatives $(\partial/\partial x_{k})f$,
$k = 1, 2, \ldots, n$, which exist in the distribution sense, are $L^{p}$
functions if and only if
$\|\tau_{h}f - f\|_{p} = (\int_{E_{n}} |f(x - h) - f(x)|^{p} \, dx)^{1/p} = O(|h|)$.
When $p = 1$, the first partial derivatives exist and are finite Borel
measures if and only if $\|\tau_{h}f - f\|_{1} = O(|h|)$.

\refstepcounter{equation}\label{eq:c1-4-6}\textbf{\theequation.}~The
translation operators can be defined on the space of finite Borel measures by
letting $\tau_{h}\mu$, where $\mu$ is in this space, be the measure whose
value at a Borel set $F$ is
$\mu(F - h) = \mu(\{x \in E_{n}; x + h \in F\})$. The measure $\mu$ is
absolutely continuous (with respect to Lebesgue measure) if and only if the
total variation of $\tau_{h}\mu - \mu$ tends to $0$ as $|h| \to 0$ (see
Bochner~\cite{Bochner4}).

\refstepcounter{equation}\label{eq:c1-4-7}\textbf{\theequation.}~If $\mu$ is a
discrete finite Borel measure then
$\mu = \sum_{k=1}^{\infty} a_{k}\delta_{x_{k}}$, where $\delta_{x_{k}}$
denotes the Dirac $\delta$ measure concentrated at the point $x_{k}$. Thus, it
follows from the definition of the Fourier transform of measures (or
distributions) that
$\hat{\mu}(t) = \sum_{k=1}^{\infty} a_{k}e^{-2\pi i t \cdot x_{k}}$. In this
case it is true that
\[
\textup{(i)}\qquad
\lim_{R\to\infty}\frac{1}{\Omega_{n}R^{n}}\int_{|t|\leq R}|\hat{\mu}(t)|^{2}\,dt
=\sum_{k=1}^{\infty}|a_{k}|^{2},
\]
where $\Omega_{n}$ denotes the volume of the unit sphere of $E_{n}$. More
generally, it can be shown that if $\mu$ is a general finite Borel measure and
$x_{1}, x_{2}, \ldots, x_{k}, \ldots$ are the distinct points of $E_{n}$ with
nonzero measures $a_{1}, a_{2}, \ldots, a_{k}, \ldots$, then equality (i) still
holds (see Wiener~\cite{Wiener1} for the one dimensional result).

\refstepcounter{equation}\label{eq:c1-4-8}\textbf{\theequation.}~Fourier
transforms of nonnegative finite Borel measures can be characterized in terms
of positive definite functions\index{positive definite function@positive definite function}. These are continuous functions, $f$, defined on
$E_{n}$ having the property
$\sum_{j=1}^{k} \sum_{i=1}^{k} f(x_{i} - x_{j})\xi_{i}\bar{\xi}_{j} \geq 0$
for any set of complex numbers $\{\xi_{1}, \xi_{2}, \ldots, \xi_{k}\}$ and
points $\{x_{1}, x_{2}, \ldots, x_{k}\}$ in $E_{n}$.
Bochner~\cite{Bochner3} has shown that $f$ is the Fourier transform of a
nonnegative finite Borel measure if and only if $f$ is positive definite (see
Rudin~\cite{Rudin1} for the extension of this result to locally compact
Abelian groups).

\refstepcounter{equation}\label{eq:c1-4-9}\textbf{\theequation.}~If there
existed constants $c_{n}, \lambda_{n}$, $n = 1, 2, \ldots, m$, such that
$e^{-\delta} = \sum_{n=1}^{m} c_{n}e^{-(\lambda_{n}\delta)^{2}}$ this would
obviously imply that if an integral is Gauss summable then it is Abel
summable. Equality (i) in the proof of Theorem~\ref{thm:c1-1-14} almost
asserts this. In view of this it is not surprising that it can be shown that
if $f$ is such that both the Abel and Gauss means, $A_{\varepsilon}(f)$ and
$G_{\varepsilon}(f)$, exist and $\lim_{\varepsilon\to0} G_{\varepsilon}(f) = l$
then $\lim_{\varepsilon\to0} A_{\varepsilon}(f)$ exists and equals $l$ (see
Cartwright~\cite{Cartwright1}, Bochner and
Chandrasekharan~\cite{BochnerChandrasekharan1}, Chapter I, Section 14).

\refstepcounter{equation}\label{eq:c1-4-10}\textbf{\theequation.}~The Poisson
integral of a function $f \in L^{p}(E_{n})$, $1 \leq p \leq \infty$, converges
to $f(x)$ for every $x$ satisfying \eqref{eq:c1-1-23} (for a proof of this
fact see G. Weiss~\cite{Weiss2}). A more general statement that will be needed
later is as follows. Suppose $\varphi(x)$ is radial (i.e.,
$\varphi(x) = \varphi_{0}(|x|)$. Assume $\varphi_{0}(r)$ is nonnegative and
decreasing in $r$ and $\int_{E_{n}} \varphi(x) \, dx = 1$. Then
$(f * \varphi_{\varepsilon})(x) \to f(x)$ for every $x$ satisfying
\eqref{eq:c1-1-23}, as $\varepsilon \to 0$. This applies in particular to the
Gauss--Weierstrass kernels). The same argument can be used to show that if
$\nu$ is a singular finite Borel measure then the Poisson--Stieltjes integral
$v(x, \varepsilon) = \int_{E_{n}} \varphi_{\varepsilon}(x - t) \, d\nu(t)$
tends to $0$ as $\varepsilon \to 0$ for almost every $x$ (this argument makes
use of the fact that, for almost every $x \in E_{n}$,
$|S_{x}|^{-1} \times \int_{S_{x}} d\nu \to 0$ as the volume, $|S_{x}|$, of a
sphere centered at $x$ tends to zero---see \eqref{frp:c2-5-6} of Chapter II or
Saks~\cite{Saks1}, Chapter IV). If $\mu$ is an arbitrary finite Borel measure
we can decompose it into its singular part $\mu_{1}$ and its absolutely
continuous part $\mu_{2}$. If $f$ is the (Radon--Nikodym) derivative of
$\mu_{2}$ we then have
$u(x, \varepsilon) = \int_{E_{n}} P(x - t, \varepsilon) \, d\mu(t) = \int_{E_{n}} P(x - t, \varepsilon) \, d\mu_{1}(t) + \int_{E_{n}} P(x - t, \varepsilon)f(t) \, dt = u_{1}(x, \varepsilon) + u_{2}(x, \varepsilon)$.
Thus, $\lim_{\varepsilon\to0} u(x, \varepsilon) = 0 + f(x) = f(x)$, for almost
every $x \in E_{n}$.

\refstepcounter{equation}\label{eq:c1-4-11}\textbf{\theequation.}~The
following additional facts concerning operators $B$ which commute with
translations are worth recording.

\textup{(a)} The only interesting cases can occur when $p \leq q$. Since if
$p > q$, then $B$ maps all elements of $\mathscr{S}$ to the zero element of
$L^{q}(E_{n})$.

\textup{(b)} If $1 \leq p \leq \infty$, and $\|B(\varphi)\|_{p} \leq M \|\varphi\|_{p}$,
for all $\varphi \in \mathscr{S}$ (that is $B\varphi = u * \varphi$, with
$B \in (L^{p}, L^{p})$), then $B \in (L^{q}, L^{q})$, with
$\left|\frac{1}{q} - \frac{1}{2}\right| \leq \left|\frac{1}{p} - \frac{1}{2}\right|$,
and $\|B\varphi\|_{q} \leq M \|\varphi\|_{q}$.

To prove (a) notice that if $\varphi \in \mathscr{D}$, and
$f_{h}(x) = \varphi(x + h) + \varphi(x)$, then
$\|f_{h}\|_{p} = 2^{1/p}\|\varphi\|_{p}$, for $h$ sufficiently large. However,
if $F_{h} = B(f_{h})$, and $\Phi = B(\varphi)$, then
$\|F_{h}\|_{q} \to 2^{1/q}\|\Phi\|_{q}$, as $|h| \to \infty$.

To prove (b) one applies the duality theorem (Theorem~\ref{thm:c1-3-20}) and
then the Riesz interpolation theorem (Theorem \ref{thm:c5-1-3} of Chapter V).

\refstepcounter{equation}\label{eq:c1-4-12}\textbf{\theequation.}~Theorem~\ref{thm:c1-3-16}
deals with bounded linear operators on $L^{p}$ into $L^{q}$ that commute with
translations. There are many linear operators mapping $\mathscr{S}$ into
$\mathscr{S}$ that commute with translations but do not have bounded linear
extensions that map some $L^{p}$ space into an $L^{q}$ space. The operator
mapping $\varphi \in \mathscr{S}$ onto its partial derivative
$\partial\varphi/\partial x_{k}$ provides an example. It follows from
Theorem~\ref{thm:c1-1-8} that this operator, $B$, applied to $\varphi$
satisfies $(B\varphi)^{\wedge}(x) = 2\pi i x_{k}\hat{\varphi}(x)$. In view of
Theorem~\ref{thm:c1-3-16} it is natural to inquire whether there exists a
tempered distribution $u$ such that $B\varphi = u * \varphi$. It is easy to
check that the distribution mapping $\varphi$ into
$(-\partial\varphi/\partial x_{k})(0)$ has this property:
$(u * \varphi)(x) = u(\tau_{x}\tilde{\varphi}) = (-\partial\tau_{x}\tilde{\varphi}/\partial x_{k})(0) = \partial\varphi/\partial x_{k}$.
Moreover, we can easily verify that $\hat{u}$ is the function $2\pi i x_{k}$.
More generally, by iteration we obtain the following result: If $P(D)$ is a
differential polynomial and $B\varphi = P(D)\varphi$ for
$\varphi \in \mathscr{S}$, then there exists $u \in \mathscr{S}'$ such that
$B\varphi = u * \varphi$, $\hat{u}(x) = P(2\pi ix)$ and
$(B\varphi)^{\wedge}(x) = P(2\pi ix)\hat{\varphi}(x) = \hat{u}(x)\hat{\varphi}(x)$.

\refstepcounter{equation}\label{eq:c1-4-13}\textbf{\theequation.}~For any
$p > 2$, there exists an $f \in L^{p}$, whose Fourier transform, as a tempered
distribution, is not a function. If the opposite were true, the closed graph
theorem would easily show that the following inequality would have to be valid
for all $f \in L^{p}$, $p > 2$:
$\int_{|x| \leq 1} |\hat{f}(x)| \, dx \leq A\|f\|_{p}$. To contradict this
inequality we take, when $n = 1$, $\hat{f}(x) = e^{-\pi(1 + i\delta)x^{2}}$,
$f(x) = e^{-\pi x^{2}/(1 + i\delta)}/(1 + i\delta)^{1/2}$, which we obtain from
Theorem~\ref{thm:c1-1-13} by an easy argument of analytic continuation. But
$\int_{|x| \leq 1} |\hat{f}(x)| \, dx = A_{1}$, and
$\|f\|_{p} \leq A_{2}\delta^{1/p - 1/2}$, as $\delta \to \infty$, which is
impossible if $p > 2$. A similar argument clearly holds for any $n$.

\refstepcounter{equation}\label{eq:c1-4-14}\textbf{\theequation.}~For further
reference we state the integral identities involving the passage to polar
coordinates and its relation with integration on the group of rotations
$SO(n)$.

\textup{(a)} Suppose $f$ is an integrable function on $E_{n}$, then
$\int_{E_{n}} f(x) \, dx = \int_{\Sigma_{n-1}} dx' (\int_{0}^{\infty} f(rx') r^{n-1} \, dr)$,
where $x'$ ranges over $\Sigma_{n-1}$ and $dx'$ is the induced Lebesgue
measure on $\Sigma_{n-1}$.

\textup{(b)} Suppose $\varphi$ is integrable on $\Sigma_{n-1}$. Then
$(1/\omega_{n-1})\int_{\Sigma_{n-1}} \varphi(x') \, dx' = \int_{SO(n)} \varphi(\sigma\mathbf{I}) \, d\sigma$,
where $d\sigma$ is Haar measure on $SO(n)$ normalized by the condition
$\int_{SO(n)} d\sigma = 1$ and $\mathbf{I} = (1, 0, \ldots, 0)$. To prove (b),
note that since $dx'$ is invariant under rotations $\sigma \in SO(n)$,
\[
\begin{aligned}(1/\omega_{n-1})\int_{\Sigma_{n-1}} \varphi(x') \, dx'&=(1/\omega_{n-1})\int_{SO(n)}\int_{\Sigma_{n-1}} \varphi(\sigma x') \, dx' \, d\sigma\\
&=(1/\omega_{n-1})\int_{\Sigma_{n-1}}\int_{SO(n)} \varphi(\sigma\sigma_{0}\mathbf{I}) \, d\sigma \, dx'\\
&=(1/\omega_{n-1})\int_{\Sigma_{n-1}}\int_{SO(n)} \varphi(\sigma\mathbf{I}) \, d\sigma \, dx'\\
&=\int_{SO(n)} \varphi(\sigma\mathbf{I}) \, d\sigma.\end{aligned}
\]
Here $\sigma_{0}$ is any rotation mapping $\mathbf{I}$ to $x'$. The next to
last equality follows by the right invariance of Haar measure $d\sigma$. For
the basic properties of Haar measure see, e.g., Weil~\cite{Weil1}.

\refstepcounter{equation}\label{eq:c1-4-15}\textbf{\theequation.}~As a
relatively straightforward consequence of Theorem~\ref{thm:c1-3-11}, one can
prove the following. Let $u$ be a tempered distribution. Then for an
appropriate integer $m$, and for each $n$-tuple $\alpha$ ($|\alpha| \leq m$),
there exist constants $a_{\alpha}$ and slowly increasing continuous functions
$f_{\alpha}$, so that
\[
u=\sum_{|\alpha|\leq m}a_{\alpha}D^{\alpha}(L_{f_{\alpha}}).
\]
Since the converse is a consequence of the material in Section 3, we see that
this kind of representation characterizes the class of tempered distributions.

\bibnoteshead

Some important treatises on classical Harmonic Analysis are
Bochner~\cite{Bochner3}, Wiener~\cite{Wiener1}, Zygmund~\cite{Zygmund1},
Titchmarsh~\cite{Titchmarsh2}, Bochner and
Chandrasekharan~\cite{BochnerChandrasekharan1}, Bari~\cite{Bari1}. For an
introductory treatment see R. Goldberg~\cite{Goldberg1} and
G. Weiss~\cite{Weiss3}. Much of the elementary analysis used in the first
chapter, in particular the introduction of polar coordinates, can be found in
Fleming~\cite{Fleming1}; in this connection also see
Blumenson~\cite{Blumenson1}. A general setting for most of the material in the
second section and the first part of the first section is the theory of locally
compact Abelian groups, for which see Weil~\cite{Weil1}, Rudin~\cite{Rudin1},
Hewitt and Ross~\cite{HewittRoss1}. We have made no attempt to discuss the
important topics originating with Wiener's Tauberian theorem which lead to the
more extensive study of the $L^{1}$ algebra. For a treatment of these topics in
the context of general locally compact Abelian groups the reader is referred
to Naimark~\cite{Naimark1} or Rudin~\cite{Rudin1}. Our proof of
Theorem~\ref{thm:c1-1-14} follows that in Bochner and
Chandrasekharan~\cite{BochnerChandrasekharan1}. The technique of the
approximation to the identity (in particular, variants of
Theorems~\ref{thm:c1-1-18} and \ref{thm:c1-1-25}) has a long history for which
we can attempt no adequate description. The formulation of
Theorem~\ref{thm:c1-1-25} presented here is based on a suggestion of R. Latzer;
a similar result can be found in Calder\'on and
Zygmund~\cite{CalderonZygmund3}, Chapter II. For the results in Functional
Analysis that were used in this chapter, in particular the weak* compactness
of the unit sphere used in the proof of Theorem~\ref{thm:c1-3-19}, we refer the
reader to Royden~\cite{Royden1}.

For the general theory of distributions and their applications see
L. Schwartz~\cite{Schwartz1}, Gelfand and Shilov~\cite{GelfandShilov1}, and
Ehrenpreis~\cite{Ehrenpreis1}. There, the reader will also find the proofs
that we omitted and the references to the classical literature of generalized
functions. For rapid introduction to the theory, in connection with the Fourier
transform, see H\"ormander~\cite{Hormander2}, Chapter I, and
Yosida~\cite{Yosida1}, Chapter VI. Theorems~\ref{thm:c1-3-16},
\ref{thm:c1-3-18}, \ref{thm:c1-3-19}, \ref{thm:c1-3-20}, dealing with operators
of type $(L^{p}, L^{q})$, have been known in various forms for some time. See,
for example, Bochner~\cite{Bochner9}, Hille and
Phillips~\cite{HillePhillips1}, Bochner and
Chandrasekharan~\cite{BochnerChandrasekharan1}, Zygmund~\cite{Zygmund1},
Chapter IV, H\"ormander~\cite{Hormander1}, where, in addition, other
references may be found.
